% ================================================================
% O singură sursă, două PDF-uri (compilează fiecare de două ori):
%
%   La clasă (interactiv, soluția se dezvăluie pe bucăți):
%     pdflatex -jobname=lesson-interactive lesson.tex
%
%   Compact (toți pașii într-un singur slide, pentru elevi / README):
%     pdflatex -jobname=lesson "\PassOptionsToClass{handout}{beamer}% ================================================================
% O singură sursă, două PDF-uri (compilează fiecare de două ori):
%
%   La clasă (interactiv, soluția se dezvăluie pe bucăți):
%     pdflatex -jobname=lesson-interactive lesson.tex
%
%   Compact (toți pașii într-un singur slide, pentru elevi / README):
%     pdflatex -jobname=lesson "\PassOptionsToClass{handout}{beamer}% ================================================================
% O singură sursă, două PDF-uri (compilează fiecare de două ori):
%
%   La clasă (interactiv, soluția se dezvăluie pe bucăți):
%     pdflatex -jobname=lesson-interactive lesson.tex
%
%   Compact (toți pașii într-un singur slide, pentru elevi / README):
%     pdflatex -jobname=lesson "\PassOptionsToClass{handout}{beamer}% ================================================================
% O singură sursă, două PDF-uri (compilează fiecare de două ori):
%
%   La clasă (interactiv, soluția se dezvăluie pe bucăți):
%     pdflatex -jobname=lesson-interactive lesson.tex
%
%   Compact (toți pașii într-un singur slide, pentru elevi / README):
%     pdflatex -jobname=lesson "\PassOptionsToClass{handout}{beamer}\input{lesson.tex}"
% ================================================================
\documentclass[9pt]{beamer}
\usepackage{../../../resources/themes/algolymp-dense}

% Mod interactiv: întrebarea apare prima, apoi soluția se dezvăluie pe bucăți
% (\item<2->, \uncover<3->{...}, \begin{remarca}<4->{Concluzie}); hint-ul se
% folosește doar unde elevii chiar trebuie împinși să se gândească.
% Totul se afișează cumulativ: ce apare la pasul 2 rămâne vizibil după.
\newenvironment<>{intrebare}{%
  \begin{uncoverenv}#1%
  \begin{block}{Întrebare}}{\end{block}\end{uncoverenv}}
% Hint și remarcă: text inline cu etichetă emerald, fără cutie
\newenvironment<>{hint}{%
  \begin{uncoverenv}#1%
  \par\textcolor{advAccent}{\textbf{Hint:}}\ }{\par\end{uncoverenv}}
\newenvironment<>{remarca}[1]{%
  \begin{uncoverenv}#2%
  \par\smallskip\textcolor{advAccent}{\textbf{#1:}}\ }{\par\end{uncoverenv}}

% Celulă de tabel vizibilă de la pasul #1
\newcommand{\at}[2]{\uncover<#1->{#2}}

\title{Sisteme de numerație \textbar\ Operații pe biți}
\subtitle{}
\author{Andrei Mihai Ciobanu}
\institute{Algolymp Summer School}
\date{}
\titlegraphic{\includegraphics[width=3cm]{../../../resources/logos/algolymp.png}}
\setbeamertemplate{frame footer}{Algolymp Summer School}

% ================================================================
\begin{document}

\begin{frame}\titlepage\end{frame}

\begin{frame}{Cuprins}\tableofcontents\end{frame}

% ================================================================
\section{Baze de numerație}

\begin{frame}[fragile]{Baze de numerație}
  Un număr poate fi scris în orice \textbf{bază} $b$, folosind cifre de la $0$ la $b-1$:
  \smallskip
  \begin{center}
  \begin{tabular}{llll}
    \toprule
    Baza & Cifre & Exemplu & Sumă de puteri \\
    \midrule
    10 (zecimal) & $0 \ldots 9$             & $425_{10}$       & $4 \cdot 10^2 + 2 \cdot 10 + 5$ \\
    2 (binar)    & $0,\ 1$                  & $1101_2 = 13$    & $2^3 + 2^2 + 2^0$ \\
    16 (hexa)    & $0\ldots 9,\ A\ldots F$  & $\text{FF}_{16} = 255$ & $15 \cdot 16 + 15$ \\
    \bottomrule
  \end{tabular}
  \end{center}
  \begin{itemize}
    \item În hexazecimal: $A=10,\ B=11,\ C=12,\ D=13,\ E=14,\ F=15$
    \item În C++: \lstinline|0b1101| = binar, \lstinline|0xFF| = hexazecimal
  \end{itemize}
  \begin{intrebare}<2->
    Care este valoarea lui $2\text{A}_{16}$ în baza $10$?
  \end{intrebare}
  \uncover<3->{$2\text{A}_{16} = 2 \cdot 16 + 10 = 42$}
\end{frame}

% ================================================================
\section{Reprezentarea numerelor în calculator}

\begin{frame}[fragile]{Reprezentarea numerelor pe biți}
  Calculatoarele lucrează doar cu \textbf{biți} \textemdash{} cifre care pot fi doar $0$ sau $1$.
  \begin{intrebare}
    Cum se poate reprezenta numărul $13$ folosind doar biți?
  \end{intrebare}
  \uncover<2->{
    Orice număr întreg se reprezintă ca o \textbf{sumă de puteri ale lui 2}:
    \[
      13 = 2^3 + 2^2 + 2^0 = 8 + 4 + 1
    \]
  }
  \begin{block}<3->{Reprezentare pe 8 biți (exemplu: 13)}
    \begin{center}
    \begin{tabular}{|c|c|c|c|c|c|c|c|}
      \hline
      $2^7$ & $2^6$ & $2^5$ & $2^4$ & $2^3$ & $2^2$ & $2^1$ & $2^0$ \\
      \hline
      0 & 0 & 0 & 0 & 1 & 1 & 0 & 1 \\
      \hline
    \end{tabular}
    \end{center}
  \end{block}
  \begin{itemize}
    \item<4-> \lstinline|int| (32 biți) $\rightarrow$ valori de la $-2^{31}$ la $2^{31}-1 \approx 2 \cdot 10^9$
    \item<4-> \lstinline|long long| (64 biți) $\rightarrow$ valori de la $-2^{63}$ la $2^{63}-1 \approx 9 \cdot 10^{18}$
  \end{itemize}
\end{frame}

% ================================================================
\section{Conversii între baze}

\begin{frame}{Conversia din baza 10 în baza $b$}
  \begin{intrebare}
    Cum se obține reprezentarea lui $13_{10}$ în baza $2$?
  \end{intrebare}
  \begin{hint}<2->
    Se împarte repetat numărul la $2$ și se rețin resturile.
  \end{hint}
  \begin{itemize}
    \item<3-> $13 \div 2 = 6$ rest $\mathbf{1}$
    \item<3-> $6 \div 2 = 3$ rest $\mathbf{0}$
    \item<3-> $3 \div 2 = 1$ rest $\mathbf{1}$
    \item<3-> $1 \div 2 = 0$ rest $\mathbf{1}$
    \item<3-> $\Rightarrow\ 13_{10} = 1101_2$
  \end{itemize}
  \begin{remarca}<4->{Algoritmul}
    Se împarte repetat la $b$, iar resturile se citesc \textbf{de jos în sus}.
  \end{remarca}
\end{frame}

\begin{frame}[fragile]{Conversia din baza 10 în baza $b$ \textemdash{} implementare}
  \begin{lstlisting}
// Scrie in digits cifrele lui n in baza b, de la cel mai
// semnificativ, si returneaza numarul de cifre
int toBase(int n, int b, int digits[]) {
    int cnt = 0;
    do {
        digits[cnt++] = n % b;
        n /= b;
    } while (n > 0);
    reverse(digits, digits + cnt); // inverseaza ordinea elementelor din intervalul [0, cnt)
    return cnt;
}
  \end{lstlisting}
  \begin{itemize}
    \item Resturile se obțin de la cifra cea mai puțin semnificativă, de aceea cifrele se inversează
    \item Bucla \lstinline|do ... while| asigură cifra $0$ pentru $n = 0$
    \item Tabloul \lstinline|digits| trebuie să aibă dimensiune suficientă
  \end{itemize}
  \beginnote{reverse \textemdash{} funcție din STL}
\end{frame}

\begin{frame}[fragile]{Conversia din baza $b$ în baza 10}
  \begin{intrebare}
    Cum se calculează valoarea lui $1101_2$ dacă se prelucrează cifrele de la stânga la dreapta?
  \end{intrebare}
  \uncover<2->{
    \[
      1101_2 = ((1 \cdot 2 + 1) \cdot 2 + 0) \cdot 2 + 1 = 13
    \]
  }
  \uncover<3->{Pas cu pas:}
  \begin{itemize}
    \item<3-> $result = 1$
    \item<3-> $result = 1 \cdot 2 + 1 = 3$
    \item<4-> $result = 3 \cdot 2 + 0 = 6$
    \item<4-> $result = 6 \cdot 2 + 1 = 13$ \checkmark
  \end{itemize}
  \begin{uncoverenv}<5->
    \begin{lstlisting}
// digits[0..cnt-1]: cifrele in baza b, de la cel mai semnificativ
int fromBase(int digits[], int cnt, int b) {
    int result = 0;
    for (int i = 0; i < cnt; i++) {
        result = result * b + digits[i];
    }
    return result;
}
    \end{lstlisting}
    La fiecare pas, rezultatul parțial se înmulțește cu $b$ și se adaugă următoarea cifră.
  \end{uncoverenv}
\end{frame}

% ================================================================
\section{Operații pe biți}

\begin{frame}{Noțiuni de bază}
  Biții unui număr se numerotează \textbf{de la dreapta la stânga}, începând cu $0$.
  \begin{intrebare}
    Care sunt biții setați ai lui $x = 89 = 1011001_2$ și ce valoare are fiecare?
  \end{intrebare}
  \uncover<2->{
    \begin{center}
    \begin{tabular}{cccccccc}
      \toprule
      Indice bit & 6 & 5 & 4 & 3 & 2 & 1 & 0 \\
      \midrule
      Bit        & 1 & 0 & 1 & 1 & 0 & 0 & 1 \\
      \midrule
      Valoare    & 64 & 32 & 16 & 8 & 4 & 2 & 1 \\
      \bottomrule
    \end{tabular}
    \end{center}
  }
  \uncover<3->{
    Biții setați: 6, 4, 3, 0 \quad$\Rightarrow$\quad $2^6 + 2^4 + 2^3 + 2^0 = 64 + 16 + 8 + 1 = 89$
  }
  \begin{itemize}
    \item<4-> \textbf{LSB} (\textit{Least Significant Bit}) \textemdash{} bitul cel mai puțin semnificativ:
      bitul $0$, cel mai din dreapta
    \item<4-> \textbf{MSB} (\textit{Most Significant Bit}) \textemdash{} bitul cel mai semnificativ:
      bitul cu cel mai mare indice setat
  \end{itemize}
\end{frame}

\begin{frame}[fragile]{NOT \textemdash{} \texttt{\textasciitilde}}
  Inversează \textbf{fiecare bit}: $0$ devine $1$, $1$ devine $0$.
  \begin{center}
  \begin{tabular}{ccc}
    \toprule
    $p$ & 0 & 1 \\
    $\sim p$ & 1 & 0 \\
    \bottomrule
  \end{tabular}
  \end{center}
  \begin{intrebare}
    Care este $\sim a$ pentru $a = 10 = 1010_2$?
  \end{intrebare}
  \uncover<2->{
    \begin{center}
    \begin{tabular}{rcccc}
      $a$:     & 1 & 0 & 1 & 0 \\
      \hline
      $\sim a$: & 0 & 1 & 0 & 1 \\
    \end{tabular}
    \end{center}
  }
  \begin{uncoverenv}<3->
    \begin{lstlisting}
int a = 0b1010; // 10
int r = ~a;     // -11
    \end{lstlisting}
    Pe 32 de biți: $\sim a = -(a+1)$, deci $\sim 10 = -11$.
  \end{uncoverenv}
  \beginnote{Complement față de 2 \textemdash{} modul în care calculatoarele stochează numerele negative}
\end{frame}

\begin{frame}[fragile]{AND \textemdash{} \texttt{\&}}
  Rezultat $1$ dacă \textbf{ambii} biți corespunzători sunt $1$.

  {\centering
  \begin{tabular}{cc|c}
    \toprule
    $p$ & $q$ & $p \mathbin{\&} q$ \\
    \midrule
    0 & 0 & 0 \\
    0 & 1 & 0 \\
    1 & 0 & 0 \\
    1 & 1 & 1 \\
    \bottomrule
  \end{tabular}\par}
  \begin{intrebare}<2->
    Care este rezultatul $a \mathbin{\&} b$ pentru $a = 1010_2$, $b = 1100_2$?
  \end{intrebare}
  \uncover<2->{\centering
  \begin{tabular}{rcccc}
    $a$: & 1 & 0 & 1 & 0 \\
    $b$: & 1 & 1 & 0 & 0 \\
    \hline
    \at{3}{$a \mathbin{\&} b$:} & \at{3}{1} & \at{3}{0} & \at{3}{0} & \at{3}{0} \\
  \end{tabular}\par}
  \begin{uncoverenv}<4->
    \begin{lstlisting}
int a = 0b1010, b = 0b1100;
int r = a & b; // 0b1000 = 8
    \end{lstlisting}
  \end{uncoverenv}
\end{frame}

\begin{frame}[fragile]{OR \textemdash{} \texttt{|}}
  Rezultat $1$ dacă \textbf{cel puțin unul} dintre biți este $1$.

  {\centering
  \begin{tabular}{cc|c}
    \toprule
    $p$ & $q$ & $p \mid q$ \\
    \midrule
    0 & 0 & 0 \\
    0 & 1 & 1 \\
    1 & 0 & 1 \\
    1 & 1 & 1 \\
    \bottomrule
  \end{tabular}\par}
  \begin{intrebare}<2->
    Care este rezultatul $a \mid b$ pentru $a = 1010_2$, $b = 1100_2$?
  \end{intrebare}
  \uncover<2->{\centering
  \begin{tabular}{rcccc}
    $a$: & 1 & 0 & 1 & 0 \\
    $b$: & 1 & 1 & 0 & 0 \\
    \hline
    \at{3}{$a \mid b$:} & \at{3}{1} & \at{3}{1} & \at{3}{1} & \at{3}{0} \\
  \end{tabular}\par}
  \begin{uncoverenv}<4->
    \begin{lstlisting}
int a = 0b1010, b = 0b1100;
int r = a | b; // 0b1110 = 14
    \end{lstlisting}
  \end{uncoverenv}
\end{frame}

\begin{frame}[fragile]{XOR \textemdash{} \texttt{\textasciicircum}}
  Rezultat $1$ dacă biții sunt \textbf{diferiți}.

  {\centering
  \begin{tabular}{cc|c}
    \toprule
    $p$ & $q$ & $p \oplus q$ \\
    \midrule
    0 & 0 & 0 \\
    0 & 1 & 1 \\
    1 & 0 & 1 \\
    1 & 1 & 0 \\
    \bottomrule
  \end{tabular}\par}
  \smallskip
  Proprietăți XOR: $a \oplus a = 0$; \quad $a \oplus 0 = a$; \quad XOR este comutativ și asociativ.
  \begin{intrebare}<2->
    Care este rezultatul $a \oplus b$ pentru $a = 1010_2$, $b = 1100_2$?
  \end{intrebare}
  \uncover<2->{\centering
  \begin{tabular}{rcccc}
    $a$: & 1 & 0 & 1 & 0 \\
    $b$: & 1 & 1 & 0 & 0 \\
    \hline
    \at{3}{$a \oplus b$:} & \at{3}{0} & \at{3}{1} & \at{3}{1} & \at{3}{0} \\
  \end{tabular}\par}
  \begin{uncoverenv}<4->
    \begin{lstlisting}
int a = 0b1010, b = 0b1100;
int r = a ^ b; // 0b0110 = 6
    \end{lstlisting}
  \end{uncoverenv}
\end{frame}

\begin{frame}[fragile]{Deplasări pe biți \textemdash{} left shift și right shift}
  \begin{intrebare}
    Ce efect au deplasările la stânga și la dreapta asupra valorii unui număr?
  \end{intrebare}
  \begin{itemize}
    \item<2-> \lstinline|x << k| $= x \cdot 2^k$ \quad (deplasare la stânga)
    \item<2-> \lstinline|x >> k| $= \lfloor x / 2^k \rfloor$ \quad (deplasare la dreapta)
  \end{itemize}
  \begin{uncoverenv}<3->
    \begin{lstlisting}
int x = 1;      // 0b00000001
int a = x << 3; // 0b00001000 = 8   (x * 2^3)
int b = x >> 0; // 0b00000001 = 1

int y = 12;     // 0b00001100
int c = y << 1; // 0b00011000 = 24  (y * 2)
int d = y >> 1; // 0b00000110 = 6   (y / 2)
int e = y >> 3; // 0b00000001 = 1   (y / 8)
    \end{lstlisting}
  \end{uncoverenv}
  \begin{remarca}<4->{Atenție}
    \lstinline|1 << k| dă overflow pentru $k \geq 31$ pe \lstinline|int|; pentru
    \lstinline|long long| se folosește \lstinline|1LL << k|.
  \end{remarca}
\end{frame}

% ================================================================
\section{Aplicații}

\begin{frame}[fragile]{Verificarea unui bit}
  \begin{intrebare}
    Cum se verifică dacă al $k$-lea bit al lui $x$ este setat?
  \end{intrebare}
  \begin{hint}<2->
    Se aduce bitul $k$ pe poziția $0$ prin deplasare, apoi se izolează cu o mască.
  \end{hint}
  \begin{uncoverenv}<3->
    \begin{lstlisting}
// Verificam daca al k-lea bit este setat
bool isSet(int x, int k) {
    return (x >> k) & 1;
}

// Caz special: x impar <=> bitul 0 setat
bool isOdd(int x) {
    return x & 1;
}
    \end{lstlisting}
  \end{uncoverenv}
  \begin{itemize}
    \item<4-> Exemplu: $x = 13 = 1101_2$, $k = 2$: $x \gg 2 = 0011_2$, iar $0011_2 \mathbin{\&} 1 = 1$ $\Rightarrow$ bitul $2$ este setat
    \item<4-> Pentru $k = 1$: $x \gg 1 = 0110_2$, iar $0110_2 \mathbin{\&} 1 = 0$ $\Rightarrow$ bitul $1$ nu este setat
  \end{itemize}
\end{frame}

\begin{frame}[fragile]{Setarea, resetarea și inversarea unui bit}
  \begin{intrebare}
    Cum se setează, se resetează și se inversează al $k$-lea bit al lui $x$?
  \end{intrebare}
  \begin{itemize}
    \item<2-> Masca \lstinline|1 << k| are setat doar bitul $k$
    \item<2-> \lstinline+|+ setează, \lstinline|& ~| resetează, \lstinline|^| inversează bitul selectat de mască
  \end{itemize}
  \begin{uncoverenv}<3->
    \begin{lstlisting}
// Setam al k-lea bit (il facem 1)
int setBit(int x, int k) {
    return x | (1 << k);
}

// Resetam al k-lea bit (il facem 0)
int clearBit(int x, int k) {
    return x & ~(1 << k);
}

// Flippuim al k-lea bit (0->1, 1->0)
int flipBit(int x, int k) {
    return x ^ (1 << k);
}
    \end{lstlisting}
  \end{uncoverenv}
\end{frame}

\begin{frame}[fragile]{Ultimii $k$ biți și puterile lui 2}
  \begin{lstlisting}
// Ultimii k biti ai lui x  (ex: k=3 => masca 0b0111)
int lastKBits(int x, int k) {
    return x & ((1 << k) - 1);
}
  \end{lstlisting}
  \begin{intrebare}
    Cum se verifică dacă $n$ este putere a lui $2$?
  \end{intrebare}
  \begin{hint}<2->
    Ce efect are $n \mathbin{\&} (n-1)$ asupra celui mai din dreapta bit de $1$?
  \end{hint}
  \begin{itemize}
    \item<3-> O putere a lui $2$ are exact un singur bit de $1$
    \item<3-> $n \mathbin{\&} (n-1)$ elimină cel mai din dreapta bit de $1$
    \item<3-> Exemplu: $4 = 100_2$,\ $3 = 011_2$,\ $4 \mathbin{\&} 3 = 0$ \checkmark
  \end{itemize}
  \begin{uncoverenv}<4->
    \begin{lstlisting}
bool isPow2(int n) {
    return n > 0 && (n & (n - 1)) == 0;
}
    \end{lstlisting}
  \end{uncoverenv}
\end{frame}

% ================================================================
\section{Funcții utile}

\begin{frame}[fragile]{Popcount \textemdash{} numărul de biți de 1}
  \begin{intrebare}
    Cum se determină numărul de biți de $1$ ai unui număr?
  \end{intrebare}
  \begin{uncoverenv}<2->
    \begin{lstlisting}
// Implementare manuala: O(nr. biti)
int popcount(int x) {
    int cnt = 0;
    while (x > 0) { cnt += x & 1; x >>= 1; }
    return cnt;
}
    \end{lstlisting}
  \end{uncoverenv}
  \begin{uncoverenv}<3->
    \begin{lstlisting}
// Built-in GCC (O(1), foarte rapid)
int cnt = __builtin_popcount(x);    // pentru int
int cnt = __builtin_popcountll(x);  // pentru long long
    \end{lstlisting}
  \end{uncoverenv}
  \begin{itemize}
    \item<4-> Exemplu: \lstinline|popcount(13)| = \lstinline|popcount(0b1101)| = 3
    \item<4-> Built-in-urile \lstinline|__builtin_*| sunt specifice GCC \textemdash{} funcționează
      pe majoritatea compilatoarelor de concurs (g++), dar nu sunt garantate pe MSVC
  \end{itemize}
\end{frame}

\begin{frame}[fragile]{Leading zeros și trailing zeros}
  \begin{itemize}
    \item \textbf{Leading zeros}: biți de $0$ \textbf{de la stânga} până la primul bit de $1$
    \item \textbf{Trailing zeros}: biți de $0$ \textbf{de la dreapta} până la primul bit de $1$
  \end{itemize}
  \begin{intrebare}
    Câți biți de $0$ sunt la stânga, respectiv la dreapta primului bit de $1$ pentru
    $x = 12 = 00001100_2$ (pe $32$ de biți)?
  \end{intrebare}
  \uncover<2->{Leading zeros $= 28$ (primii $28$ de biți sunt $0$), trailing zeros $= 2$.}
  \begin{uncoverenv}<3->
    \begin{lstlisting}
int lz = __builtin_clz(x);  // leading zeros  (int)
int tz = __builtin_ctz(x);  // trailing zeros (int)
// Variante pentru long long: __builtin_clzll, __builtin_ctzll

// Atentie: __builtin_clz(0) si __builtin_ctz(0) sunt nedefinite!
    \end{lstlisting}
  \end{uncoverenv}
  \begin{itemize}
    \item<4-> Trailing zeros al lui $x$ = indicele celui mai mic bit setat al lui $x$
    \item<4-> Exemplu: \lstinline|__builtin_ctz(12)| = \lstinline|__builtin_ctz(0b1100)| = 2
  \end{itemize}
\end{frame}

\begin{frame}[fragile]{Lowest set bit \textemdash{} \lstinline|x & (-x)|}
  \begin{intrebare}
    Cum se izolează cel mai din dreapta bit de $1$ (\textbf{lowest set bit}) al lui $x$?
  \end{intrebare}
  \begin{hint}<2->
    Reprezentarea \textbf{complement față de 2}: $-x = \,\sim\! x + 1$ (se inversează toți biții și se adaugă $1$).
  \end{hint}
  \begin{uncoverenv}<3->
    \begin{lstlisting}
int lowestBit(int x) { return x & (-x); }
// x = 12 = 0b1100  =>  lowestBit(12) = 0b0100 = 4
    \end{lstlisting}
  \end{uncoverenv}
  \begin{itemize}
    \item<4-> La adunarea lui $1$, carry-ul se propagă exact până la primul bit de $1$ din $x$
    \item<4-> Astfel $-x$ coincide cu $x$ începând de la primul bit de $1$ (inclusiv), iar restul este inversat
    \item<4-> $x \mathbin{\&} (-x)$ păstrează doar acel bit comun
  \end{itemize}
  \uncover<5->{
    Exemplu: $x = 12 = \mathbf{1100}_2$,\quad $-x = \mathbf{0100}_2$ (în 4 biți),\quad
    $x \mathbin{\&} (-x) = 0100_2 = 4$
  }
\end{frame}

% ================================================================
\section{Recapitulare}

\begin{frame}[fragile]{Recapitulare \textemdash{} Operații uzuale}
  \begin{center}
  \begin{tabular}{ll}
    \toprule
    Operație & Cod C++ \\
    \midrule
    Bitul $k$ e setat?   & \lstinline|(x >> k) & 1| \\
    $x$ e impar?         & \lstinline|x & 1| \\
    Setează bitul $k$    & \lstinline+x | (1 << k)+ \\
    Resetează bitul $k$  & \lstinline|x & ~(1 << k)| \\
    Flip bitul $k$       & \lstinline|x ^ (1 << k)| \\
    Ultimii $k$ biți     & \lstinline|x & ((1 << k) - 1)| \\
    $n$ e putere de 2?   & \lstinline|n > 0 && !(n & (n-1))| \\
    Lowest set bit       & \lstinline|x & (-x)| \\
    Număr de biți de 1   & \lstinline|__builtin_popcount(x)| \\
    Trailing zeros       & \lstinline|__builtin_ctz(x)| \\
    \bottomrule
  \end{tabular}
  \end{center}
\end{frame}

\begin{frame}{Probleme de antrenament}
  \begin{block}{Ușoare}
    \begin{enumerate}
      \item \href{https://www.pbinfo.ro/probleme/2577/getbit}{\textbf{getbit}}
            \textcolor{gray}{(pbinfo.ro)}
      \item \href{https://www.pbinfo.ro/probleme/2580/setbit0}{\textbf{setbit0}}
            \textcolor{gray}{(pbinfo.ro)}
      \item \href{https://www.pbinfo.ro/probleme/2581/setbit1}{\textbf{setbit1}}
            \textcolor{gray}{(pbinfo.ro)}
      \item \href{https://www.pbinfo.ro/probleme/2711/puterebiti}{\textbf{puterebiti}}
            \textcolor{gray}{(pbinfo.ro)}
    \end{enumerate}
  \end{block}
  \begin{block}{Medii}
    \begin{enumerate}
      \item \href{https://code.algolymp.com/problems/4271}{\textbf{Base Bit Report}}
            \textcolor{gray}{(AlgoCode \#4271)}
      \item \href{https://www.pbinfo.ro/probleme/2799/addone}{\textbf{addone}}
            \textcolor{gray}{(pbinfo.ro)}
      \item \href{https://www.pbinfo.ro/probleme/2624/invbits}{\textbf{invbits}}
            \textcolor{gray}{(pbinfo.ro)}
    \end{enumerate}
  \end{block}
\end{frame}

\begin{frame}{Resurse}
  \begin{enumerate}
    \item \href{https://www.pbinfo.ro/articole/5562/baze-de-numeratie}{\textbf{pbinfo.ro}}
          \textemdash{} Baze de numerație (RO)
    \item \href{https://www.pbinfo.ro/articole/62/operatii-pe-biti}{\textbf{pbinfo.ro}}
          \textemdash{} Operații pe biți (RO, ușor de înțeles)
    \item \href{https://cp-algorithms.com/algebra/bit-manipulation.html}{\textbf{cp-algorithms.com}}
          \textemdash{} Bit manipulation (EN)
    \item \href{https://cppi.sync.ro/materia/lectii/baze\%20de\%20numeratie,\%20scrierea\%20constantelor\%20si\%20operatii\%20pe\%20biti.pdf}{\textbf{cppi.sync.ro}}
          \textemdash{} Material mai stufos de Marius Nicoli (RO)
    \item \href{https://en.wikipedia.org/wiki/Two\%27s_complement}{\textbf{wikipedia.org}}
          \textemdash{} Two's complement (EN)
  \end{enumerate}
\end{frame}

\end{document}
"
% ================================================================
\documentclass[9pt]{beamer}
\usepackage{../../../resources/themes/algolymp-dense}

% Mod interactiv: întrebarea apare prima, apoi soluția se dezvăluie pe bucăți
% (\item<2->, \uncover<3->{...}, \begin{remarca}<4->{Concluzie}); hint-ul se
% folosește doar unde elevii chiar trebuie împinși să se gândească.
% Totul se afișează cumulativ: ce apare la pasul 2 rămâne vizibil după.
\newenvironment<>{intrebare}{%
  \begin{uncoverenv}#1%
  \begin{block}{Întrebare}}{\end{block}\end{uncoverenv}}
% Hint și remarcă: text inline cu etichetă emerald, fără cutie
\newenvironment<>{hint}{%
  \begin{uncoverenv}#1%
  \par\textcolor{advAccent}{\textbf{Hint:}}\ }{\par\end{uncoverenv}}
\newenvironment<>{remarca}[1]{%
  \begin{uncoverenv}#2%
  \par\smallskip\textcolor{advAccent}{\textbf{#1:}}\ }{\par\end{uncoverenv}}

% Celulă de tabel vizibilă de la pasul #1
\newcommand{\at}[2]{\uncover<#1->{#2}}

\title{Sisteme de numerație \textbar\ Operații pe biți}
\subtitle{}
\author{Andrei Mihai Ciobanu}
\institute{Algolymp Summer School}
\date{}
\titlegraphic{\includegraphics[width=3cm]{../../../resources/logos/algolymp.png}}
\setbeamertemplate{frame footer}{Algolymp Summer School}

% ================================================================
\begin{document}

\begin{frame}\titlepage\end{frame}

\begin{frame}{Cuprins}\tableofcontents\end{frame}

% ================================================================
\section{Baze de numerație}

\begin{frame}[fragile]{Baze de numerație}
  Un număr poate fi scris în orice \textbf{bază} $b$, folosind cifre de la $0$ la $b-1$:
  \smallskip
  \begin{center}
  \begin{tabular}{llll}
    \toprule
    Baza & Cifre & Exemplu & Sumă de puteri \\
    \midrule
    10 (zecimal) & $0 \ldots 9$             & $425_{10}$       & $4 \cdot 10^2 + 2 \cdot 10 + 5$ \\
    2 (binar)    & $0,\ 1$                  & $1101_2 = 13$    & $2^3 + 2^2 + 2^0$ \\
    16 (hexa)    & $0\ldots 9,\ A\ldots F$  & $\text{FF}_{16} = 255$ & $15 \cdot 16 + 15$ \\
    \bottomrule
  \end{tabular}
  \end{center}
  \begin{itemize}
    \item În hexazecimal: $A=10,\ B=11,\ C=12,\ D=13,\ E=14,\ F=15$
    \item În C++: \lstinline|0b1101| = binar, \lstinline|0xFF| = hexazecimal
  \end{itemize}
  \begin{intrebare}<2->
    Care este valoarea lui $2\text{A}_{16}$ în baza $10$?
  \end{intrebare}
  \uncover<3->{$2\text{A}_{16} = 2 \cdot 16 + 10 = 42$}
\end{frame}

% ================================================================
\section{Reprezentarea numerelor în calculator}

\begin{frame}[fragile]{Reprezentarea numerelor pe biți}
  Calculatoarele lucrează doar cu \textbf{biți} \textemdash{} cifre care pot fi doar $0$ sau $1$.
  \begin{intrebare}
    Cum se poate reprezenta numărul $13$ folosind doar biți?
  \end{intrebare}
  \uncover<2->{
    Orice număr întreg se reprezintă ca o \textbf{sumă de puteri ale lui 2}:
    \[
      13 = 2^3 + 2^2 + 2^0 = 8 + 4 + 1
    \]
  }
  \begin{block}<3->{Reprezentare pe 8 biți (exemplu: 13)}
    \begin{center}
    \begin{tabular}{|c|c|c|c|c|c|c|c|}
      \hline
      $2^7$ & $2^6$ & $2^5$ & $2^4$ & $2^3$ & $2^2$ & $2^1$ & $2^0$ \\
      \hline
      0 & 0 & 0 & 0 & 1 & 1 & 0 & 1 \\
      \hline
    \end{tabular}
    \end{center}
  \end{block}
  \begin{itemize}
    \item<4-> \lstinline|int| (32 biți) $\rightarrow$ valori de la $-2^{31}$ la $2^{31}-1 \approx 2 \cdot 10^9$
    \item<4-> \lstinline|long long| (64 biți) $\rightarrow$ valori de la $-2^{63}$ la $2^{63}-1 \approx 9 \cdot 10^{18}$
  \end{itemize}
\end{frame}

% ================================================================
\section{Conversii între baze}

\begin{frame}{Conversia din baza 10 în baza $b$}
  \begin{intrebare}
    Cum se obține reprezentarea lui $13_{10}$ în baza $2$?
  \end{intrebare}
  \begin{hint}<2->
    Se împarte repetat numărul la $2$ și se rețin resturile.
  \end{hint}
  \begin{itemize}
    \item<3-> $13 \div 2 = 6$ rest $\mathbf{1}$
    \item<3-> $6 \div 2 = 3$ rest $\mathbf{0}$
    \item<3-> $3 \div 2 = 1$ rest $\mathbf{1}$
    \item<3-> $1 \div 2 = 0$ rest $\mathbf{1}$
    \item<3-> $\Rightarrow\ 13_{10} = 1101_2$
  \end{itemize}
  \begin{remarca}<4->{Algoritmul}
    Se împarte repetat la $b$, iar resturile se citesc \textbf{de jos în sus}.
  \end{remarca}
\end{frame}

\begin{frame}[fragile]{Conversia din baza 10 în baza $b$ \textemdash{} implementare}
  \begin{lstlisting}
// Scrie in digits cifrele lui n in baza b, de la cel mai
// semnificativ, si returneaza numarul de cifre
int toBase(int n, int b, int digits[]) {
    int cnt = 0;
    do {
        digits[cnt++] = n % b;
        n /= b;
    } while (n > 0);
    reverse(digits, digits + cnt); // inverseaza ordinea elementelor din intervalul [0, cnt)
    return cnt;
}
  \end{lstlisting}
  \begin{itemize}
    \item Resturile se obțin de la cifra cea mai puțin semnificativă, de aceea cifrele se inversează
    \item Bucla \lstinline|do ... while| asigură cifra $0$ pentru $n = 0$
    \item Tabloul \lstinline|digits| trebuie să aibă dimensiune suficientă
  \end{itemize}
  \beginnote{reverse \textemdash{} funcție din STL}
\end{frame}

\begin{frame}[fragile]{Conversia din baza $b$ în baza 10}
  \begin{intrebare}
    Cum se calculează valoarea lui $1101_2$ dacă se prelucrează cifrele de la stânga la dreapta?
  \end{intrebare}
  \uncover<2->{
    \[
      1101_2 = ((1 \cdot 2 + 1) \cdot 2 + 0) \cdot 2 + 1 = 13
    \]
  }
  \uncover<3->{Pas cu pas:}
  \begin{itemize}
    \item<3-> $result = 1$
    \item<3-> $result = 1 \cdot 2 + 1 = 3$
    \item<4-> $result = 3 \cdot 2 + 0 = 6$
    \item<4-> $result = 6 \cdot 2 + 1 = 13$ \checkmark
  \end{itemize}
  \begin{uncoverenv}<5->
    \begin{lstlisting}
// digits[0..cnt-1]: cifrele in baza b, de la cel mai semnificativ
int fromBase(int digits[], int cnt, int b) {
    int result = 0;
    for (int i = 0; i < cnt; i++) {
        result = result * b + digits[i];
    }
    return result;
}
    \end{lstlisting}
    La fiecare pas, rezultatul parțial se înmulțește cu $b$ și se adaugă următoarea cifră.
  \end{uncoverenv}
\end{frame}

% ================================================================
\section{Operații pe biți}

\begin{frame}{Noțiuni de bază}
  Biții unui număr se numerotează \textbf{de la dreapta la stânga}, începând cu $0$.
  \begin{intrebare}
    Care sunt biții setați ai lui $x = 89 = 1011001_2$ și ce valoare are fiecare?
  \end{intrebare}
  \uncover<2->{
    \begin{center}
    \begin{tabular}{cccccccc}
      \toprule
      Indice bit & 6 & 5 & 4 & 3 & 2 & 1 & 0 \\
      \midrule
      Bit        & 1 & 0 & 1 & 1 & 0 & 0 & 1 \\
      \midrule
      Valoare    & 64 & 32 & 16 & 8 & 4 & 2 & 1 \\
      \bottomrule
    \end{tabular}
    \end{center}
  }
  \uncover<3->{
    Biții setați: 6, 4, 3, 0 \quad$\Rightarrow$\quad $2^6 + 2^4 + 2^3 + 2^0 = 64 + 16 + 8 + 1 = 89$
  }
  \begin{itemize}
    \item<4-> \textbf{LSB} (\textit{Least Significant Bit}) \textemdash{} bitul cel mai puțin semnificativ:
      bitul $0$, cel mai din dreapta
    \item<4-> \textbf{MSB} (\textit{Most Significant Bit}) \textemdash{} bitul cel mai semnificativ:
      bitul cu cel mai mare indice setat
  \end{itemize}
\end{frame}

\begin{frame}[fragile]{NOT \textemdash{} \texttt{\textasciitilde}}
  Inversează \textbf{fiecare bit}: $0$ devine $1$, $1$ devine $0$.
  \begin{center}
  \begin{tabular}{ccc}
    \toprule
    $p$ & 0 & 1 \\
    $\sim p$ & 1 & 0 \\
    \bottomrule
  \end{tabular}
  \end{center}
  \begin{intrebare}
    Care este $\sim a$ pentru $a = 10 = 1010_2$?
  \end{intrebare}
  \uncover<2->{
    \begin{center}
    \begin{tabular}{rcccc}
      $a$:     & 1 & 0 & 1 & 0 \\
      \hline
      $\sim a$: & 0 & 1 & 0 & 1 \\
    \end{tabular}
    \end{center}
  }
  \begin{uncoverenv}<3->
    \begin{lstlisting}
int a = 0b1010; // 10
int r = ~a;     // -11
    \end{lstlisting}
    Pe 32 de biți: $\sim a = -(a+1)$, deci $\sim 10 = -11$.
  \end{uncoverenv}
  \beginnote{Complement față de 2 \textemdash{} modul în care calculatoarele stochează numerele negative}
\end{frame}

\begin{frame}[fragile]{AND \textemdash{} \texttt{\&}}
  Rezultat $1$ dacă \textbf{ambii} biți corespunzători sunt $1$.

  {\centering
  \begin{tabular}{cc|c}
    \toprule
    $p$ & $q$ & $p \mathbin{\&} q$ \\
    \midrule
    0 & 0 & 0 \\
    0 & 1 & 0 \\
    1 & 0 & 0 \\
    1 & 1 & 1 \\
    \bottomrule
  \end{tabular}\par}
  \begin{intrebare}<2->
    Care este rezultatul $a \mathbin{\&} b$ pentru $a = 1010_2$, $b = 1100_2$?
  \end{intrebare}
  \uncover<2->{\centering
  \begin{tabular}{rcccc}
    $a$: & 1 & 0 & 1 & 0 \\
    $b$: & 1 & 1 & 0 & 0 \\
    \hline
    \at{3}{$a \mathbin{\&} b$:} & \at{3}{1} & \at{3}{0} & \at{3}{0} & \at{3}{0} \\
  \end{tabular}\par}
  \begin{uncoverenv}<4->
    \begin{lstlisting}
int a = 0b1010, b = 0b1100;
int r = a & b; // 0b1000 = 8
    \end{lstlisting}
  \end{uncoverenv}
\end{frame}

\begin{frame}[fragile]{OR \textemdash{} \texttt{|}}
  Rezultat $1$ dacă \textbf{cel puțin unul} dintre biți este $1$.

  {\centering
  \begin{tabular}{cc|c}
    \toprule
    $p$ & $q$ & $p \mid q$ \\
    \midrule
    0 & 0 & 0 \\
    0 & 1 & 1 \\
    1 & 0 & 1 \\
    1 & 1 & 1 \\
    \bottomrule
  \end{tabular}\par}
  \begin{intrebare}<2->
    Care este rezultatul $a \mid b$ pentru $a = 1010_2$, $b = 1100_2$?
  \end{intrebare}
  \uncover<2->{\centering
  \begin{tabular}{rcccc}
    $a$: & 1 & 0 & 1 & 0 \\
    $b$: & 1 & 1 & 0 & 0 \\
    \hline
    \at{3}{$a \mid b$:} & \at{3}{1} & \at{3}{1} & \at{3}{1} & \at{3}{0} \\
  \end{tabular}\par}
  \begin{uncoverenv}<4->
    \begin{lstlisting}
int a = 0b1010, b = 0b1100;
int r = a | b; // 0b1110 = 14
    \end{lstlisting}
  \end{uncoverenv}
\end{frame}

\begin{frame}[fragile]{XOR \textemdash{} \texttt{\textasciicircum}}
  Rezultat $1$ dacă biții sunt \textbf{diferiți}.

  {\centering
  \begin{tabular}{cc|c}
    \toprule
    $p$ & $q$ & $p \oplus q$ \\
    \midrule
    0 & 0 & 0 \\
    0 & 1 & 1 \\
    1 & 0 & 1 \\
    1 & 1 & 0 \\
    \bottomrule
  \end{tabular}\par}
  \smallskip
  Proprietăți XOR: $a \oplus a = 0$; \quad $a \oplus 0 = a$; \quad XOR este comutativ și asociativ.
  \begin{intrebare}<2->
    Care este rezultatul $a \oplus b$ pentru $a = 1010_2$, $b = 1100_2$?
  \end{intrebare}
  \uncover<2->{\centering
  \begin{tabular}{rcccc}
    $a$: & 1 & 0 & 1 & 0 \\
    $b$: & 1 & 1 & 0 & 0 \\
    \hline
    \at{3}{$a \oplus b$:} & \at{3}{0} & \at{3}{1} & \at{3}{1} & \at{3}{0} \\
  \end{tabular}\par}
  \begin{uncoverenv}<4->
    \begin{lstlisting}
int a = 0b1010, b = 0b1100;
int r = a ^ b; // 0b0110 = 6
    \end{lstlisting}
  \end{uncoverenv}
\end{frame}

\begin{frame}[fragile]{Deplasări pe biți \textemdash{} left shift și right shift}
  \begin{intrebare}
    Ce efect au deplasările la stânga și la dreapta asupra valorii unui număr?
  \end{intrebare}
  \begin{itemize}
    \item<2-> \lstinline|x << k| $= x \cdot 2^k$ \quad (deplasare la stânga)
    \item<2-> \lstinline|x >> k| $= \lfloor x / 2^k \rfloor$ \quad (deplasare la dreapta)
  \end{itemize}
  \begin{uncoverenv}<3->
    \begin{lstlisting}
int x = 1;      // 0b00000001
int a = x << 3; // 0b00001000 = 8   (x * 2^3)
int b = x >> 0; // 0b00000001 = 1

int y = 12;     // 0b00001100
int c = y << 1; // 0b00011000 = 24  (y * 2)
int d = y >> 1; // 0b00000110 = 6   (y / 2)
int e = y >> 3; // 0b00000001 = 1   (y / 8)
    \end{lstlisting}
  \end{uncoverenv}
  \begin{remarca}<4->{Atenție}
    \lstinline|1 << k| dă overflow pentru $k \geq 31$ pe \lstinline|int|; pentru
    \lstinline|long long| se folosește \lstinline|1LL << k|.
  \end{remarca}
\end{frame}

% ================================================================
\section{Aplicații}

\begin{frame}[fragile]{Verificarea unui bit}
  \begin{intrebare}
    Cum se verifică dacă al $k$-lea bit al lui $x$ este setat?
  \end{intrebare}
  \begin{hint}<2->
    Se aduce bitul $k$ pe poziția $0$ prin deplasare, apoi se izolează cu o mască.
  \end{hint}
  \begin{uncoverenv}<3->
    \begin{lstlisting}
// Verificam daca al k-lea bit este setat
bool isSet(int x, int k) {
    return (x >> k) & 1;
}

// Caz special: x impar <=> bitul 0 setat
bool isOdd(int x) {
    return x & 1;
}
    \end{lstlisting}
  \end{uncoverenv}
  \begin{itemize}
    \item<4-> Exemplu: $x = 13 = 1101_2$, $k = 2$: $x \gg 2 = 0011_2$, iar $0011_2 \mathbin{\&} 1 = 1$ $\Rightarrow$ bitul $2$ este setat
    \item<4-> Pentru $k = 1$: $x \gg 1 = 0110_2$, iar $0110_2 \mathbin{\&} 1 = 0$ $\Rightarrow$ bitul $1$ nu este setat
  \end{itemize}
\end{frame}

\begin{frame}[fragile]{Setarea, resetarea și inversarea unui bit}
  \begin{intrebare}
    Cum se setează, se resetează și se inversează al $k$-lea bit al lui $x$?
  \end{intrebare}
  \begin{itemize}
    \item<2-> Masca \lstinline|1 << k| are setat doar bitul $k$
    \item<2-> \lstinline+|+ setează, \lstinline|& ~| resetează, \lstinline|^| inversează bitul selectat de mască
  \end{itemize}
  \begin{uncoverenv}<3->
    \begin{lstlisting}
// Setam al k-lea bit (il facem 1)
int setBit(int x, int k) {
    return x | (1 << k);
}

// Resetam al k-lea bit (il facem 0)
int clearBit(int x, int k) {
    return x & ~(1 << k);
}

// Flippuim al k-lea bit (0->1, 1->0)
int flipBit(int x, int k) {
    return x ^ (1 << k);
}
    \end{lstlisting}
  \end{uncoverenv}
\end{frame}

\begin{frame}[fragile]{Ultimii $k$ biți și puterile lui 2}
  \begin{lstlisting}
// Ultimii k biti ai lui x  (ex: k=3 => masca 0b0111)
int lastKBits(int x, int k) {
    return x & ((1 << k) - 1);
}
  \end{lstlisting}
  \begin{intrebare}
    Cum se verifică dacă $n$ este putere a lui $2$?
  \end{intrebare}
  \begin{hint}<2->
    Ce efect are $n \mathbin{\&} (n-1)$ asupra celui mai din dreapta bit de $1$?
  \end{hint}
  \begin{itemize}
    \item<3-> O putere a lui $2$ are exact un singur bit de $1$
    \item<3-> $n \mathbin{\&} (n-1)$ elimină cel mai din dreapta bit de $1$
    \item<3-> Exemplu: $4 = 100_2$,\ $3 = 011_2$,\ $4 \mathbin{\&} 3 = 0$ \checkmark
  \end{itemize}
  \begin{uncoverenv}<4->
    \begin{lstlisting}
bool isPow2(int n) {
    return n > 0 && (n & (n - 1)) == 0;
}
    \end{lstlisting}
  \end{uncoverenv}
\end{frame}

% ================================================================
\section{Funcții utile}

\begin{frame}[fragile]{Popcount \textemdash{} numărul de biți de 1}
  \begin{intrebare}
    Cum se determină numărul de biți de $1$ ai unui număr?
  \end{intrebare}
  \begin{uncoverenv}<2->
    \begin{lstlisting}
// Implementare manuala: O(nr. biti)
int popcount(int x) {
    int cnt = 0;
    while (x > 0) { cnt += x & 1; x >>= 1; }
    return cnt;
}
    \end{lstlisting}
  \end{uncoverenv}
  \begin{uncoverenv}<3->
    \begin{lstlisting}
// Built-in GCC (O(1), foarte rapid)
int cnt = __builtin_popcount(x);    // pentru int
int cnt = __builtin_popcountll(x);  // pentru long long
    \end{lstlisting}
  \end{uncoverenv}
  \begin{itemize}
    \item<4-> Exemplu: \lstinline|popcount(13)| = \lstinline|popcount(0b1101)| = 3
    \item<4-> Built-in-urile \lstinline|__builtin_*| sunt specifice GCC \textemdash{} funcționează
      pe majoritatea compilatoarelor de concurs (g++), dar nu sunt garantate pe MSVC
  \end{itemize}
\end{frame}

\begin{frame}[fragile]{Leading zeros și trailing zeros}
  \begin{itemize}
    \item \textbf{Leading zeros}: biți de $0$ \textbf{de la stânga} până la primul bit de $1$
    \item \textbf{Trailing zeros}: biți de $0$ \textbf{de la dreapta} până la primul bit de $1$
  \end{itemize}
  \begin{intrebare}
    Câți biți de $0$ sunt la stânga, respectiv la dreapta primului bit de $1$ pentru
    $x = 12 = 00001100_2$ (pe $32$ de biți)?
  \end{intrebare}
  \uncover<2->{Leading zeros $= 28$ (primii $28$ de biți sunt $0$), trailing zeros $= 2$.}
  \begin{uncoverenv}<3->
    \begin{lstlisting}
int lz = __builtin_clz(x);  // leading zeros  (int)
int tz = __builtin_ctz(x);  // trailing zeros (int)
// Variante pentru long long: __builtin_clzll, __builtin_ctzll

// Atentie: __builtin_clz(0) si __builtin_ctz(0) sunt nedefinite!
    \end{lstlisting}
  \end{uncoverenv}
  \begin{itemize}
    \item<4-> Trailing zeros al lui $x$ = indicele celui mai mic bit setat al lui $x$
    \item<4-> Exemplu: \lstinline|__builtin_ctz(12)| = \lstinline|__builtin_ctz(0b1100)| = 2
  \end{itemize}
\end{frame}

\begin{frame}[fragile]{Lowest set bit \textemdash{} \lstinline|x & (-x)|}
  \begin{intrebare}
    Cum se izolează cel mai din dreapta bit de $1$ (\textbf{lowest set bit}) al lui $x$?
  \end{intrebare}
  \begin{hint}<2->
    Reprezentarea \textbf{complement față de 2}: $-x = \,\sim\! x + 1$ (se inversează toți biții și se adaugă $1$).
  \end{hint}
  \begin{uncoverenv}<3->
    \begin{lstlisting}
int lowestBit(int x) { return x & (-x); }
// x = 12 = 0b1100  =>  lowestBit(12) = 0b0100 = 4
    \end{lstlisting}
  \end{uncoverenv}
  \begin{itemize}
    \item<4-> La adunarea lui $1$, carry-ul se propagă exact până la primul bit de $1$ din $x$
    \item<4-> Astfel $-x$ coincide cu $x$ începând de la primul bit de $1$ (inclusiv), iar restul este inversat
    \item<4-> $x \mathbin{\&} (-x)$ păstrează doar acel bit comun
  \end{itemize}
  \uncover<5->{
    Exemplu: $x = 12 = \mathbf{1100}_2$,\quad $-x = \mathbf{0100}_2$ (în 4 biți),\quad
    $x \mathbin{\&} (-x) = 0100_2 = 4$
  }
\end{frame}

% ================================================================
\section{Recapitulare}

\begin{frame}[fragile]{Recapitulare \textemdash{} Operații uzuale}
  \begin{center}
  \begin{tabular}{ll}
    \toprule
    Operație & Cod C++ \\
    \midrule
    Bitul $k$ e setat?   & \lstinline|(x >> k) & 1| \\
    $x$ e impar?         & \lstinline|x & 1| \\
    Setează bitul $k$    & \lstinline+x | (1 << k)+ \\
    Resetează bitul $k$  & \lstinline|x & ~(1 << k)| \\
    Flip bitul $k$       & \lstinline|x ^ (1 << k)| \\
    Ultimii $k$ biți     & \lstinline|x & ((1 << k) - 1)| \\
    $n$ e putere de 2?   & \lstinline|n > 0 && !(n & (n-1))| \\
    Lowest set bit       & \lstinline|x & (-x)| \\
    Număr de biți de 1   & \lstinline|__builtin_popcount(x)| \\
    Trailing zeros       & \lstinline|__builtin_ctz(x)| \\
    \bottomrule
  \end{tabular}
  \end{center}
\end{frame}

\begin{frame}{Probleme de antrenament}
  \begin{block}{Ușoare}
    \begin{enumerate}
      \item \href{https://www.pbinfo.ro/probleme/2577/getbit}{\textbf{getbit}}
            \textcolor{gray}{(pbinfo.ro)}
      \item \href{https://www.pbinfo.ro/probleme/2580/setbit0}{\textbf{setbit0}}
            \textcolor{gray}{(pbinfo.ro)}
      \item \href{https://www.pbinfo.ro/probleme/2581/setbit1}{\textbf{setbit1}}
            \textcolor{gray}{(pbinfo.ro)}
      \item \href{https://www.pbinfo.ro/probleme/2711/puterebiti}{\textbf{puterebiti}}
            \textcolor{gray}{(pbinfo.ro)}
    \end{enumerate}
  \end{block}
  \begin{block}{Medii}
    \begin{enumerate}
      \item \href{https://code.algolymp.com/problems/4271}{\textbf{Base Bit Report}}
            \textcolor{gray}{(AlgoCode \#4271)}
      \item \href{https://www.pbinfo.ro/probleme/2799/addone}{\textbf{addone}}
            \textcolor{gray}{(pbinfo.ro)}
      \item \href{https://www.pbinfo.ro/probleme/2624/invbits}{\textbf{invbits}}
            \textcolor{gray}{(pbinfo.ro)}
    \end{enumerate}
  \end{block}
\end{frame}

\begin{frame}{Resurse}
  \begin{enumerate}
    \item \href{https://www.pbinfo.ro/articole/5562/baze-de-numeratie}{\textbf{pbinfo.ro}}
          \textemdash{} Baze de numerație (RO)
    \item \href{https://www.pbinfo.ro/articole/62/operatii-pe-biti}{\textbf{pbinfo.ro}}
          \textemdash{} Operații pe biți (RO, ușor de înțeles)
    \item \href{https://cp-algorithms.com/algebra/bit-manipulation.html}{\textbf{cp-algorithms.com}}
          \textemdash{} Bit manipulation (EN)
    \item \href{https://cppi.sync.ro/materia/lectii/baze\%20de\%20numeratie,\%20scrierea\%20constantelor\%20si\%20operatii\%20pe\%20biti.pdf}{\textbf{cppi.sync.ro}}
          \textemdash{} Material mai stufos de Marius Nicoli (RO)
    \item \href{https://en.wikipedia.org/wiki/Two\%27s_complement}{\textbf{wikipedia.org}}
          \textemdash{} Two's complement (EN)
  \end{enumerate}
\end{frame}

\end{document}
"
% ================================================================
\documentclass[9pt]{beamer}
\usepackage{../../../resources/themes/algolymp-dense}

% Mod interactiv: întrebarea apare prima, apoi soluția se dezvăluie pe bucăți
% (\item<2->, \uncover<3->{...}, \begin{remarca}<4->{Concluzie}); hint-ul se
% folosește doar unde elevii chiar trebuie împinși să se gândească.
% Totul se afișează cumulativ: ce apare la pasul 2 rămâne vizibil după.
\newenvironment<>{intrebare}{%
  \begin{uncoverenv}#1%
  \begin{block}{Întrebare}}{\end{block}\end{uncoverenv}}
% Hint și remarcă: text inline cu etichetă emerald, fără cutie
\newenvironment<>{hint}{%
  \begin{uncoverenv}#1%
  \par\textcolor{advAccent}{\textbf{Hint:}}\ }{\par\end{uncoverenv}}
\newenvironment<>{remarca}[1]{%
  \begin{uncoverenv}#2%
  \par\smallskip\textcolor{advAccent}{\textbf{#1:}}\ }{\par\end{uncoverenv}}

% Celulă de tabel vizibilă de la pasul #1
\newcommand{\at}[2]{\uncover<#1->{#2}}

\title{Sisteme de numerație \textbar\ Operații pe biți}
\subtitle{}
\author{Andrei Mihai Ciobanu}
\institute{Algolymp Summer School}
\date{}
\titlegraphic{\includegraphics[width=3cm]{../../../resources/logos/algolymp.png}}
\setbeamertemplate{frame footer}{Algolymp Summer School}

% ================================================================
\begin{document}

\begin{frame}\titlepage\end{frame}

\begin{frame}{Cuprins}\tableofcontents\end{frame}

% ================================================================
\section{Baze de numerație}

\begin{frame}[fragile]{Baze de numerație}
  Un număr poate fi scris în orice \textbf{bază} $b$, folosind cifre de la $0$ la $b-1$:
  \smallskip
  \begin{center}
  \begin{tabular}{llll}
    \toprule
    Baza & Cifre & Exemplu & Sumă de puteri \\
    \midrule
    10 (zecimal) & $0 \ldots 9$             & $425_{10}$       & $4 \cdot 10^2 + 2 \cdot 10 + 5$ \\
    2 (binar)    & $0,\ 1$                  & $1101_2 = 13$    & $2^3 + 2^2 + 2^0$ \\
    16 (hexa)    & $0\ldots 9,\ A\ldots F$  & $\text{FF}_{16} = 255$ & $15 \cdot 16 + 15$ \\
    \bottomrule
  \end{tabular}
  \end{center}
  \begin{itemize}
    \item În hexazecimal: $A=10,\ B=11,\ C=12,\ D=13,\ E=14,\ F=15$
    \item În C++: \lstinline|0b1101| = binar, \lstinline|0xFF| = hexazecimal
  \end{itemize}
  \begin{intrebare}<2->
    Care este valoarea lui $2\text{A}_{16}$ în baza $10$?
  \end{intrebare}
  \uncover<3->{$2\text{A}_{16} = 2 \cdot 16 + 10 = 42$}
\end{frame}

% ================================================================
\section{Reprezentarea numerelor în calculator}

\begin{frame}[fragile]{Reprezentarea numerelor pe biți}
  Calculatoarele lucrează doar cu \textbf{biți} \textemdash{} cifre care pot fi doar $0$ sau $1$.
  \begin{intrebare}
    Cum se poate reprezenta numărul $13$ folosind doar biți?
  \end{intrebare}
  \uncover<2->{
    Orice număr întreg se reprezintă ca o \textbf{sumă de puteri ale lui 2}:
    \[
      13 = 2^3 + 2^2 + 2^0 = 8 + 4 + 1
    \]
  }
  \begin{block}<3->{Reprezentare pe 8 biți (exemplu: 13)}
    \begin{center}
    \begin{tabular}{|c|c|c|c|c|c|c|c|}
      \hline
      $2^7$ & $2^6$ & $2^5$ & $2^4$ & $2^3$ & $2^2$ & $2^1$ & $2^0$ \\
      \hline
      0 & 0 & 0 & 0 & 1 & 1 & 0 & 1 \\
      \hline
    \end{tabular}
    \end{center}
  \end{block}
  \begin{itemize}
    \item<4-> \lstinline|int| (32 biți) $\rightarrow$ valori de la $-2^{31}$ la $2^{31}-1 \approx 2 \cdot 10^9$
    \item<4-> \lstinline|long long| (64 biți) $\rightarrow$ valori de la $-2^{63}$ la $2^{63}-1 \approx 9 \cdot 10^{18}$
  \end{itemize}
\end{frame}

% ================================================================
\section{Conversii între baze}

\begin{frame}{Conversia din baza 10 în baza $b$}
  \begin{intrebare}
    Cum se obține reprezentarea lui $13_{10}$ în baza $2$?
  \end{intrebare}
  \begin{hint}<2->
    Se împarte repetat numărul la $2$ și se rețin resturile.
  \end{hint}
  \begin{itemize}
    \item<3-> $13 \div 2 = 6$ rest $\mathbf{1}$
    \item<3-> $6 \div 2 = 3$ rest $\mathbf{0}$
    \item<3-> $3 \div 2 = 1$ rest $\mathbf{1}$
    \item<3-> $1 \div 2 = 0$ rest $\mathbf{1}$
    \item<3-> $\Rightarrow\ 13_{10} = 1101_2$
  \end{itemize}
  \begin{remarca}<4->{Algoritmul}
    Se împarte repetat la $b$, iar resturile se citesc \textbf{de jos în sus}.
  \end{remarca}
\end{frame}

\begin{frame}[fragile]{Conversia din baza 10 în baza $b$ \textemdash{} implementare}
  \begin{lstlisting}
// Scrie in digits cifrele lui n in baza b, de la cel mai
// semnificativ, si returneaza numarul de cifre
int toBase(int n, int b, int digits[]) {
    int cnt = 0;
    do {
        digits[cnt++] = n % b;
        n /= b;
    } while (n > 0);
    reverse(digits, digits + cnt); // inverseaza ordinea elementelor din intervalul [0, cnt)
    return cnt;
}
  \end{lstlisting}
  \begin{itemize}
    \item Resturile se obțin de la cifra cea mai puțin semnificativă, de aceea cifrele se inversează
    \item Bucla \lstinline|do ... while| asigură cifra $0$ pentru $n = 0$
    \item Tabloul \lstinline|digits| trebuie să aibă dimensiune suficientă
  \end{itemize}
  \beginnote{reverse \textemdash{} funcție din STL}
\end{frame}

\begin{frame}[fragile]{Conversia din baza $b$ în baza 10}
  \begin{intrebare}
    Cum se calculează valoarea lui $1101_2$ dacă se prelucrează cifrele de la stânga la dreapta?
  \end{intrebare}
  \uncover<2->{
    \[
      1101_2 = ((1 \cdot 2 + 1) \cdot 2 + 0) \cdot 2 + 1 = 13
    \]
  }
  \uncover<3->{Pas cu pas:}
  \begin{itemize}
    \item<3-> $result = 1$
    \item<3-> $result = 1 \cdot 2 + 1 = 3$
    \item<4-> $result = 3 \cdot 2 + 0 = 6$
    \item<4-> $result = 6 \cdot 2 + 1 = 13$ \checkmark
  \end{itemize}
  \begin{uncoverenv}<5->
    \begin{lstlisting}
// digits[0..cnt-1]: cifrele in baza b, de la cel mai semnificativ
int fromBase(int digits[], int cnt, int b) {
    int result = 0;
    for (int i = 0; i < cnt; i++) {
        result = result * b + digits[i];
    }
    return result;
}
    \end{lstlisting}
    La fiecare pas, rezultatul parțial se înmulțește cu $b$ și se adaugă următoarea cifră.
  \end{uncoverenv}
\end{frame}

% ================================================================
\section{Operații pe biți}

\begin{frame}{Noțiuni de bază}
  Biții unui număr se numerotează \textbf{de la dreapta la stânga}, începând cu $0$.
  \begin{intrebare}
    Care sunt biții setați ai lui $x = 89 = 1011001_2$ și ce valoare are fiecare?
  \end{intrebare}
  \uncover<2->{
    \begin{center}
    \begin{tabular}{cccccccc}
      \toprule
      Indice bit & 6 & 5 & 4 & 3 & 2 & 1 & 0 \\
      \midrule
      Bit        & 1 & 0 & 1 & 1 & 0 & 0 & 1 \\
      \midrule
      Valoare    & 64 & 32 & 16 & 8 & 4 & 2 & 1 \\
      \bottomrule
    \end{tabular}
    \end{center}
  }
  \uncover<3->{
    Biții setați: 6, 4, 3, 0 \quad$\Rightarrow$\quad $2^6 + 2^4 + 2^3 + 2^0 = 64 + 16 + 8 + 1 = 89$
  }
  \begin{itemize}
    \item<4-> \textbf{LSB} (\textit{Least Significant Bit}) \textemdash{} bitul cel mai puțin semnificativ:
      bitul $0$, cel mai din dreapta
    \item<4-> \textbf{MSB} (\textit{Most Significant Bit}) \textemdash{} bitul cel mai semnificativ:
      bitul cu cel mai mare indice setat
  \end{itemize}
\end{frame}

\begin{frame}[fragile]{NOT \textemdash{} \texttt{\textasciitilde}}
  Inversează \textbf{fiecare bit}: $0$ devine $1$, $1$ devine $0$.
  \begin{center}
  \begin{tabular}{ccc}
    \toprule
    $p$ & 0 & 1 \\
    $\sim p$ & 1 & 0 \\
    \bottomrule
  \end{tabular}
  \end{center}
  \begin{intrebare}
    Care este $\sim a$ pentru $a = 10 = 1010_2$?
  \end{intrebare}
  \uncover<2->{
    \begin{center}
    \begin{tabular}{rcccc}
      $a$:     & 1 & 0 & 1 & 0 \\
      \hline
      $\sim a$: & 0 & 1 & 0 & 1 \\
    \end{tabular}
    \end{center}
  }
  \begin{uncoverenv}<3->
    \begin{lstlisting}
int a = 0b1010; // 10
int r = ~a;     // -11
    \end{lstlisting}
    Pe 32 de biți: $\sim a = -(a+1)$, deci $\sim 10 = -11$.
  \end{uncoverenv}
  \beginnote{Complement față de 2 \textemdash{} modul în care calculatoarele stochează numerele negative}
\end{frame}

\begin{frame}[fragile]{AND \textemdash{} \texttt{\&}}
  Rezultat $1$ dacă \textbf{ambii} biți corespunzători sunt $1$.

  {\centering
  \begin{tabular}{cc|c}
    \toprule
    $p$ & $q$ & $p \mathbin{\&} q$ \\
    \midrule
    0 & 0 & 0 \\
    0 & 1 & 0 \\
    1 & 0 & 0 \\
    1 & 1 & 1 \\
    \bottomrule
  \end{tabular}\par}
  \begin{intrebare}<2->
    Care este rezultatul $a \mathbin{\&} b$ pentru $a = 1010_2$, $b = 1100_2$?
  \end{intrebare}
  \uncover<2->{\centering
  \begin{tabular}{rcccc}
    $a$: & 1 & 0 & 1 & 0 \\
    $b$: & 1 & 1 & 0 & 0 \\
    \hline
    \at{3}{$a \mathbin{\&} b$:} & \at{3}{1} & \at{3}{0} & \at{3}{0} & \at{3}{0} \\
  \end{tabular}\par}
  \begin{uncoverenv}<4->
    \begin{lstlisting}
int a = 0b1010, b = 0b1100;
int r = a & b; // 0b1000 = 8
    \end{lstlisting}
  \end{uncoverenv}
\end{frame}

\begin{frame}[fragile]{OR \textemdash{} \texttt{|}}
  Rezultat $1$ dacă \textbf{cel puțin unul} dintre biți este $1$.

  {\centering
  \begin{tabular}{cc|c}
    \toprule
    $p$ & $q$ & $p \mid q$ \\
    \midrule
    0 & 0 & 0 \\
    0 & 1 & 1 \\
    1 & 0 & 1 \\
    1 & 1 & 1 \\
    \bottomrule
  \end{tabular}\par}
  \begin{intrebare}<2->
    Care este rezultatul $a \mid b$ pentru $a = 1010_2$, $b = 1100_2$?
  \end{intrebare}
  \uncover<2->{\centering
  \begin{tabular}{rcccc}
    $a$: & 1 & 0 & 1 & 0 \\
    $b$: & 1 & 1 & 0 & 0 \\
    \hline
    \at{3}{$a \mid b$:} & \at{3}{1} & \at{3}{1} & \at{3}{1} & \at{3}{0} \\
  \end{tabular}\par}
  \begin{uncoverenv}<4->
    \begin{lstlisting}
int a = 0b1010, b = 0b1100;
int r = a | b; // 0b1110 = 14
    \end{lstlisting}
  \end{uncoverenv}
\end{frame}

\begin{frame}[fragile]{XOR \textemdash{} \texttt{\textasciicircum}}
  Rezultat $1$ dacă biții sunt \textbf{diferiți}.

  {\centering
  \begin{tabular}{cc|c}
    \toprule
    $p$ & $q$ & $p \oplus q$ \\
    \midrule
    0 & 0 & 0 \\
    0 & 1 & 1 \\
    1 & 0 & 1 \\
    1 & 1 & 0 \\
    \bottomrule
  \end{tabular}\par}
  \smallskip
  Proprietăți XOR: $a \oplus a = 0$; \quad $a \oplus 0 = a$; \quad XOR este comutativ și asociativ.
  \begin{intrebare}<2->
    Care este rezultatul $a \oplus b$ pentru $a = 1010_2$, $b = 1100_2$?
  \end{intrebare}
  \uncover<2->{\centering
  \begin{tabular}{rcccc}
    $a$: & 1 & 0 & 1 & 0 \\
    $b$: & 1 & 1 & 0 & 0 \\
    \hline
    \at{3}{$a \oplus b$:} & \at{3}{0} & \at{3}{1} & \at{3}{1} & \at{3}{0} \\
  \end{tabular}\par}
  \begin{uncoverenv}<4->
    \begin{lstlisting}
int a = 0b1010, b = 0b1100;
int r = a ^ b; // 0b0110 = 6
    \end{lstlisting}
  \end{uncoverenv}
\end{frame}

\begin{frame}[fragile]{Deplasări pe biți \textemdash{} left shift și right shift}
  \begin{intrebare}
    Ce efect au deplasările la stânga și la dreapta asupra valorii unui număr?
  \end{intrebare}
  \begin{itemize}
    \item<2-> \lstinline|x << k| $= x \cdot 2^k$ \quad (deplasare la stânga)
    \item<2-> \lstinline|x >> k| $= \lfloor x / 2^k \rfloor$ \quad (deplasare la dreapta)
  \end{itemize}
  \begin{uncoverenv}<3->
    \begin{lstlisting}
int x = 1;      // 0b00000001
int a = x << 3; // 0b00001000 = 8   (x * 2^3)
int b = x >> 0; // 0b00000001 = 1

int y = 12;     // 0b00001100
int c = y << 1; // 0b00011000 = 24  (y * 2)
int d = y >> 1; // 0b00000110 = 6   (y / 2)
int e = y >> 3; // 0b00000001 = 1   (y / 8)
    \end{lstlisting}
  \end{uncoverenv}
  \begin{remarca}<4->{Atenție}
    \lstinline|1 << k| dă overflow pentru $k \geq 31$ pe \lstinline|int|; pentru
    \lstinline|long long| se folosește \lstinline|1LL << k|.
  \end{remarca}
\end{frame}

% ================================================================
\section{Aplicații}

\begin{frame}[fragile]{Verificarea unui bit}
  \begin{intrebare}
    Cum se verifică dacă al $k$-lea bit al lui $x$ este setat?
  \end{intrebare}
  \begin{hint}<2->
    Se aduce bitul $k$ pe poziția $0$ prin deplasare, apoi se izolează cu o mască.
  \end{hint}
  \begin{uncoverenv}<3->
    \begin{lstlisting}
// Verificam daca al k-lea bit este setat
bool isSet(int x, int k) {
    return (x >> k) & 1;
}

// Caz special: x impar <=> bitul 0 setat
bool isOdd(int x) {
    return x & 1;
}
    \end{lstlisting}
  \end{uncoverenv}
  \begin{itemize}
    \item<4-> Exemplu: $x = 13 = 1101_2$, $k = 2$: $x \gg 2 = 0011_2$, iar $0011_2 \mathbin{\&} 1 = 1$ $\Rightarrow$ bitul $2$ este setat
    \item<4-> Pentru $k = 1$: $x \gg 1 = 0110_2$, iar $0110_2 \mathbin{\&} 1 = 0$ $\Rightarrow$ bitul $1$ nu este setat
  \end{itemize}
\end{frame}

\begin{frame}[fragile]{Setarea, resetarea și inversarea unui bit}
  \begin{intrebare}
    Cum se setează, se resetează și se inversează al $k$-lea bit al lui $x$?
  \end{intrebare}
  \begin{itemize}
    \item<2-> Masca \lstinline|1 << k| are setat doar bitul $k$
    \item<2-> \lstinline+|+ setează, \lstinline|& ~| resetează, \lstinline|^| inversează bitul selectat de mască
  \end{itemize}
  \begin{uncoverenv}<3->
    \begin{lstlisting}
// Setam al k-lea bit (il facem 1)
int setBit(int x, int k) {
    return x | (1 << k);
}

// Resetam al k-lea bit (il facem 0)
int clearBit(int x, int k) {
    return x & ~(1 << k);
}

// Flippuim al k-lea bit (0->1, 1->0)
int flipBit(int x, int k) {
    return x ^ (1 << k);
}
    \end{lstlisting}
  \end{uncoverenv}
\end{frame}

\begin{frame}[fragile]{Ultimii $k$ biți și puterile lui 2}
  \begin{lstlisting}
// Ultimii k biti ai lui x  (ex: k=3 => masca 0b0111)
int lastKBits(int x, int k) {
    return x & ((1 << k) - 1);
}
  \end{lstlisting}
  \begin{intrebare}
    Cum se verifică dacă $n$ este putere a lui $2$?
  \end{intrebare}
  \begin{hint}<2->
    Ce efect are $n \mathbin{\&} (n-1)$ asupra celui mai din dreapta bit de $1$?
  \end{hint}
  \begin{itemize}
    \item<3-> O putere a lui $2$ are exact un singur bit de $1$
    \item<3-> $n \mathbin{\&} (n-1)$ elimină cel mai din dreapta bit de $1$
    \item<3-> Exemplu: $4 = 100_2$,\ $3 = 011_2$,\ $4 \mathbin{\&} 3 = 0$ \checkmark
  \end{itemize}
  \begin{uncoverenv}<4->
    \begin{lstlisting}
bool isPow2(int n) {
    return n > 0 && (n & (n - 1)) == 0;
}
    \end{lstlisting}
  \end{uncoverenv}
\end{frame}

% ================================================================
\section{Funcții utile}

\begin{frame}[fragile]{Popcount \textemdash{} numărul de biți de 1}
  \begin{intrebare}
    Cum se determină numărul de biți de $1$ ai unui număr?
  \end{intrebare}
  \begin{uncoverenv}<2->
    \begin{lstlisting}
// Implementare manuala: O(nr. biti)
int popcount(int x) {
    int cnt = 0;
    while (x > 0) { cnt += x & 1; x >>= 1; }
    return cnt;
}
    \end{lstlisting}
  \end{uncoverenv}
  \begin{uncoverenv}<3->
    \begin{lstlisting}
// Built-in GCC (O(1), foarte rapid)
int cnt = __builtin_popcount(x);    // pentru int
int cnt = __builtin_popcountll(x);  // pentru long long
    \end{lstlisting}
  \end{uncoverenv}
  \begin{itemize}
    \item<4-> Exemplu: \lstinline|popcount(13)| = \lstinline|popcount(0b1101)| = 3
    \item<4-> Built-in-urile \lstinline|__builtin_*| sunt specifice GCC \textemdash{} funcționează
      pe majoritatea compilatoarelor de concurs (g++), dar nu sunt garantate pe MSVC
  \end{itemize}
\end{frame}

\begin{frame}[fragile]{Leading zeros și trailing zeros}
  \begin{itemize}
    \item \textbf{Leading zeros}: biți de $0$ \textbf{de la stânga} până la primul bit de $1$
    \item \textbf{Trailing zeros}: biți de $0$ \textbf{de la dreapta} până la primul bit de $1$
  \end{itemize}
  \begin{intrebare}
    Câți biți de $0$ sunt la stânga, respectiv la dreapta primului bit de $1$ pentru
    $x = 12 = 00001100_2$ (pe $32$ de biți)?
  \end{intrebare}
  \uncover<2->{Leading zeros $= 28$ (primii $28$ de biți sunt $0$), trailing zeros $= 2$.}
  \begin{uncoverenv}<3->
    \begin{lstlisting}
int lz = __builtin_clz(x);  // leading zeros  (int)
int tz = __builtin_ctz(x);  // trailing zeros (int)
// Variante pentru long long: __builtin_clzll, __builtin_ctzll

// Atentie: __builtin_clz(0) si __builtin_ctz(0) sunt nedefinite!
    \end{lstlisting}
  \end{uncoverenv}
  \begin{itemize}
    \item<4-> Trailing zeros al lui $x$ = indicele celui mai mic bit setat al lui $x$
    \item<4-> Exemplu: \lstinline|__builtin_ctz(12)| = \lstinline|__builtin_ctz(0b1100)| = 2
  \end{itemize}
\end{frame}

\begin{frame}[fragile]{Lowest set bit \textemdash{} \lstinline|x & (-x)|}
  \begin{intrebare}
    Cum se izolează cel mai din dreapta bit de $1$ (\textbf{lowest set bit}) al lui $x$?
  \end{intrebare}
  \begin{hint}<2->
    Reprezentarea \textbf{complement față de 2}: $-x = \,\sim\! x + 1$ (se inversează toți biții și se adaugă $1$).
  \end{hint}
  \begin{uncoverenv}<3->
    \begin{lstlisting}
int lowestBit(int x) { return x & (-x); }
// x = 12 = 0b1100  =>  lowestBit(12) = 0b0100 = 4
    \end{lstlisting}
  \end{uncoverenv}
  \begin{itemize}
    \item<4-> La adunarea lui $1$, carry-ul se propagă exact până la primul bit de $1$ din $x$
    \item<4-> Astfel $-x$ coincide cu $x$ începând de la primul bit de $1$ (inclusiv), iar restul este inversat
    \item<4-> $x \mathbin{\&} (-x)$ păstrează doar acel bit comun
  \end{itemize}
  \uncover<5->{
    Exemplu: $x = 12 = \mathbf{1100}_2$,\quad $-x = \mathbf{0100}_2$ (în 4 biți),\quad
    $x \mathbin{\&} (-x) = 0100_2 = 4$
  }
\end{frame}

% ================================================================
\section{Recapitulare}

\begin{frame}[fragile]{Recapitulare \textemdash{} Operații uzuale}
  \begin{center}
  \begin{tabular}{ll}
    \toprule
    Operație & Cod C++ \\
    \midrule
    Bitul $k$ e setat?   & \lstinline|(x >> k) & 1| \\
    $x$ e impar?         & \lstinline|x & 1| \\
    Setează bitul $k$    & \lstinline+x | (1 << k)+ \\
    Resetează bitul $k$  & \lstinline|x & ~(1 << k)| \\
    Flip bitul $k$       & \lstinline|x ^ (1 << k)| \\
    Ultimii $k$ biți     & \lstinline|x & ((1 << k) - 1)| \\
    $n$ e putere de 2?   & \lstinline|n > 0 && !(n & (n-1))| \\
    Lowest set bit       & \lstinline|x & (-x)| \\
    Număr de biți de 1   & \lstinline|__builtin_popcount(x)| \\
    Trailing zeros       & \lstinline|__builtin_ctz(x)| \\
    \bottomrule
  \end{tabular}
  \end{center}
\end{frame}

\begin{frame}{Probleme de antrenament}
  \begin{block}{Ușoare}
    \begin{enumerate}
      \item \href{https://www.pbinfo.ro/probleme/2577/getbit}{\textbf{getbit}}
            \textcolor{gray}{(pbinfo.ro)}
      \item \href{https://www.pbinfo.ro/probleme/2580/setbit0}{\textbf{setbit0}}
            \textcolor{gray}{(pbinfo.ro)}
      \item \href{https://www.pbinfo.ro/probleme/2581/setbit1}{\textbf{setbit1}}
            \textcolor{gray}{(pbinfo.ro)}
      \item \href{https://www.pbinfo.ro/probleme/2711/puterebiti}{\textbf{puterebiti}}
            \textcolor{gray}{(pbinfo.ro)}
    \end{enumerate}
  \end{block}
  \begin{block}{Medii}
    \begin{enumerate}
      \item \href{https://code.algolymp.com/problems/4271}{\textbf{Base Bit Report}}
            \textcolor{gray}{(AlgoCode \#4271)}
      \item \href{https://www.pbinfo.ro/probleme/2799/addone}{\textbf{addone}}
            \textcolor{gray}{(pbinfo.ro)}
      \item \href{https://www.pbinfo.ro/probleme/2624/invbits}{\textbf{invbits}}
            \textcolor{gray}{(pbinfo.ro)}
    \end{enumerate}
  \end{block}
\end{frame}

\begin{frame}{Resurse}
  \begin{enumerate}
    \item \href{https://www.pbinfo.ro/articole/5562/baze-de-numeratie}{\textbf{pbinfo.ro}}
          \textemdash{} Baze de numerație (RO)
    \item \href{https://www.pbinfo.ro/articole/62/operatii-pe-biti}{\textbf{pbinfo.ro}}
          \textemdash{} Operații pe biți (RO, ușor de înțeles)
    \item \href{https://cp-algorithms.com/algebra/bit-manipulation.html}{\textbf{cp-algorithms.com}}
          \textemdash{} Bit manipulation (EN)
    \item \href{https://cppi.sync.ro/materia/lectii/baze\%20de\%20numeratie,\%20scrierea\%20constantelor\%20si\%20operatii\%20pe\%20biti.pdf}{\textbf{cppi.sync.ro}}
          \textemdash{} Material mai stufos de Marius Nicoli (RO)
    \item \href{https://en.wikipedia.org/wiki/Two\%27s_complement}{\textbf{wikipedia.org}}
          \textemdash{} Two's complement (EN)
  \end{enumerate}
\end{frame}

\end{document}
"
% ================================================================
\documentclass[9pt]{beamer}
\usepackage{../../../resources/themes/algolymp-dense}

% Mod interactiv: întrebarea apare prima, apoi soluția se dezvăluie pe bucăți
% (\item<2->, \uncover<3->{...}, \begin{remarca}<4->{Concluzie}); hint-ul se
% folosește doar unde elevii chiar trebuie împinși să se gândească.
% Totul se afișează cumulativ: ce apare la pasul 2 rămâne vizibil după.
\newenvironment<>{intrebare}{%
  \begin{uncoverenv}#1%
  \begin{block}{Întrebare}}{\end{block}\end{uncoverenv}}
% Hint și remarcă: text inline cu etichetă emerald, fără cutie
\newenvironment<>{hint}{%
  \begin{uncoverenv}#1%
  \par\textcolor{advAccent}{\textbf{Hint:}}\ }{\par\end{uncoverenv}}
\newenvironment<>{remarca}[1]{%
  \begin{uncoverenv}#2%
  \par\smallskip\textcolor{advAccent}{\textbf{#1:}}\ }{\par\end{uncoverenv}}

% Celulă de tabel vizibilă de la pasul #1
\newcommand{\at}[2]{\uncover<#1->{#2}}

\title{Sisteme de numerație \textbar\ Operații pe biți}
\subtitle{}
\author{Andrei Mihai Ciobanu}
\institute{Algolymp Summer School}
\date{}
\titlegraphic{\includegraphics[width=3cm]{../../../resources/logos/algolymp.png}}
\setbeamertemplate{frame footer}{Algolymp Summer School}

% ================================================================
\begin{document}

\begin{frame}\titlepage\end{frame}

\begin{frame}{Cuprins}\tableofcontents\end{frame}

% ================================================================
\section{Baze de numerație}

\begin{frame}[fragile]{Baze de numerație}
  Un număr poate fi scris în orice \textbf{bază} $b$, folosind cifre de la $0$ la $b-1$:
  \smallskip
  \begin{center}
  \begin{tabular}{llll}
    \toprule
    Baza & Cifre & Exemplu & Sumă de puteri \\
    \midrule
    10 (zecimal) & $0 \ldots 9$             & $425_{10}$       & $4 \cdot 10^2 + 2 \cdot 10 + 5$ \\
    2 (binar)    & $0,\ 1$                  & $1101_2 = 13$    & $2^3 + 2^2 + 2^0$ \\
    16 (hexa)    & $0\ldots 9,\ A\ldots F$  & $\text{FF}_{16} = 255$ & $15 \cdot 16 + 15$ \\
    \bottomrule
  \end{tabular}
  \end{center}
  \begin{itemize}
    \item În hexazecimal: $A=10,\ B=11,\ C=12,\ D=13,\ E=14,\ F=15$
    \item În C++: \lstinline|0b1101| = binar, \lstinline|0xFF| = hexazecimal
  \end{itemize}
  \begin{intrebare}<2->
    Care este valoarea lui $2\text{A}_{16}$ în baza $10$?
  \end{intrebare}
  \uncover<3->{$2\text{A}_{16} = 2 \cdot 16 + 10 = 42$}
\end{frame}

% ================================================================
\section{Reprezentarea numerelor în calculator}

\begin{frame}[fragile]{Reprezentarea numerelor pe biți}
  Calculatoarele lucrează doar cu \textbf{biți} \textemdash{} cifre care pot fi doar $0$ sau $1$.
  \begin{intrebare}
    Cum se poate reprezenta numărul $13$ folosind doar biți?
  \end{intrebare}
  \uncover<2->{
    Orice număr întreg se reprezintă ca o \textbf{sumă de puteri ale lui 2}:
    \[
      13 = 2^3 + 2^2 + 2^0 = 8 + 4 + 1
    \]
  }
  \begin{block}<3->{Reprezentare pe 8 biți (exemplu: 13)}
    \begin{center}
    \begin{tabular}{|c|c|c|c|c|c|c|c|}
      \hline
      $2^7$ & $2^6$ & $2^5$ & $2^4$ & $2^3$ & $2^2$ & $2^1$ & $2^0$ \\
      \hline
      0 & 0 & 0 & 0 & 1 & 1 & 0 & 1 \\
      \hline
    \end{tabular}
    \end{center}
  \end{block}
  \begin{itemize}
    \item<4-> \lstinline|int| (32 biți) $\rightarrow$ valori de la $-2^{31}$ la $2^{31}-1 \approx 2 \cdot 10^9$
    \item<4-> \lstinline|long long| (64 biți) $\rightarrow$ valori de la $-2^{63}$ la $2^{63}-1 \approx 9 \cdot 10^{18}$
  \end{itemize}
\end{frame}

% ================================================================
\section{Conversii între baze}

\begin{frame}{Conversia din baza 10 în baza $b$}
  \begin{intrebare}
    Cum se obține reprezentarea lui $13_{10}$ în baza $2$?
  \end{intrebare}
  \begin{hint}<2->
    Se împarte repetat numărul la $2$ și se rețin resturile.
  \end{hint}
  \begin{itemize}
    \item<3-> $13 \div 2 = 6$ rest $\mathbf{1}$
    \item<3-> $6 \div 2 = 3$ rest $\mathbf{0}$
    \item<3-> $3 \div 2 = 1$ rest $\mathbf{1}$
    \item<3-> $1 \div 2 = 0$ rest $\mathbf{1}$
    \item<3-> $\Rightarrow\ 13_{10} = 1101_2$
  \end{itemize}
  \begin{remarca}<4->{Algoritmul}
    Se împarte repetat la $b$, iar resturile se citesc \textbf{de jos în sus}.
  \end{remarca}
\end{frame}

\begin{frame}[fragile]{Conversia din baza 10 în baza $b$ \textemdash{} implementare}
  \begin{lstlisting}
// Scrie in digits cifrele lui n in baza b, de la cel mai
// semnificativ, si returneaza numarul de cifre
int toBase(int n, int b, int digits[]) {
    int cnt = 0;
    do {
        digits[cnt++] = n % b;
        n /= b;
    } while (n > 0);
    reverse(digits, digits + cnt); // inverseaza ordinea elementelor din intervalul [0, cnt)
    return cnt;
}
  \end{lstlisting}
  \begin{itemize}
    \item Resturile se obțin de la cifra cea mai puțin semnificativă, de aceea cifrele se inversează
    \item Bucla \lstinline|do ... while| asigură cifra $0$ pentru $n = 0$
    \item Tabloul \lstinline|digits| trebuie să aibă dimensiune suficientă
  \end{itemize}
  \beginnote{reverse \textemdash{} funcție din STL}
\end{frame}

\begin{frame}[fragile]{Conversia din baza $b$ în baza 10}
  \begin{intrebare}
    Cum se calculează valoarea lui $1101_2$ dacă se prelucrează cifrele de la stânga la dreapta?
  \end{intrebare}
  \uncover<2->{
    \[
      1101_2 = ((1 \cdot 2 + 1) \cdot 2 + 0) \cdot 2 + 1 = 13
    \]
  }
  \uncover<3->{Pas cu pas:}
  \begin{itemize}
    \item<3-> $result = 1$
    \item<3-> $result = 1 \cdot 2 + 1 = 3$
    \item<4-> $result = 3 \cdot 2 + 0 = 6$
    \item<4-> $result = 6 \cdot 2 + 1 = 13$ \checkmark
  \end{itemize}
  \begin{uncoverenv}<5->
    \begin{lstlisting}
// digits[0..cnt-1]: cifrele in baza b, de la cel mai semnificativ
int fromBase(int digits[], int cnt, int b) {
    int result = 0;
    for (int i = 0; i < cnt; i++) {
        result = result * b + digits[i];
    }
    return result;
}
    \end{lstlisting}
    La fiecare pas, rezultatul parțial se înmulțește cu $b$ și se adaugă următoarea cifră.
  \end{uncoverenv}
\end{frame}

% ================================================================
\section{Operații pe biți}

\begin{frame}{Noțiuni de bază}
  Biții unui număr se numerotează \textbf{de la dreapta la stânga}, începând cu $0$.
  \begin{intrebare}
    Care sunt biții setați ai lui $x = 89 = 1011001_2$ și ce valoare are fiecare?
  \end{intrebare}
  \uncover<2->{
    \begin{center}
    \begin{tabular}{cccccccc}
      \toprule
      Indice bit & 6 & 5 & 4 & 3 & 2 & 1 & 0 \\
      \midrule
      Bit        & 1 & 0 & 1 & 1 & 0 & 0 & 1 \\
      \midrule
      Valoare    & 64 & 32 & 16 & 8 & 4 & 2 & 1 \\
      \bottomrule
    \end{tabular}
    \end{center}
  }
  \uncover<3->{
    Biții setați: 6, 4, 3, 0 \quad$\Rightarrow$\quad $2^6 + 2^4 + 2^3 + 2^0 = 64 + 16 + 8 + 1 = 89$
  }
  \begin{itemize}
    \item<4-> \textbf{LSB} (\textit{Least Significant Bit}) \textemdash{} bitul cel mai puțin semnificativ:
      bitul $0$, cel mai din dreapta
    \item<4-> \textbf{MSB} (\textit{Most Significant Bit}) \textemdash{} bitul cel mai semnificativ:
      bitul cu cel mai mare indice setat
  \end{itemize}
\end{frame}

\begin{frame}[fragile]{NOT \textemdash{} \texttt{\textasciitilde}}
  Inversează \textbf{fiecare bit}: $0$ devine $1$, $1$ devine $0$.
  \begin{center}
  \begin{tabular}{ccc}
    \toprule
    $p$ & 0 & 1 \\
    $\sim p$ & 1 & 0 \\
    \bottomrule
  \end{tabular}
  \end{center}
  \begin{intrebare}
    Care este $\sim a$ pentru $a = 10 = 1010_2$?
  \end{intrebare}
  \uncover<2->{
    \begin{center}
    \begin{tabular}{rcccc}
      $a$:     & 1 & 0 & 1 & 0 \\
      \hline
      $\sim a$: & 0 & 1 & 0 & 1 \\
    \end{tabular}
    \end{center}
  }
  \begin{uncoverenv}<3->
    \begin{lstlisting}
int a = 0b1010; // 10
int r = ~a;     // -11
    \end{lstlisting}
    Pe 32 de biți: $\sim a = -(a+1)$, deci $\sim 10 = -11$.
  \end{uncoverenv}
  \beginnote{Complement față de 2 \textemdash{} modul în care calculatoarele stochează numerele negative}
\end{frame}

\begin{frame}[fragile]{AND \textemdash{} \texttt{\&}}
  Rezultat $1$ dacă \textbf{ambii} biți corespunzători sunt $1$.

  {\centering
  \begin{tabular}{cc|c}
    \toprule
    $p$ & $q$ & $p \mathbin{\&} q$ \\
    \midrule
    0 & 0 & 0 \\
    0 & 1 & 0 \\
    1 & 0 & 0 \\
    1 & 1 & 1 \\
    \bottomrule
  \end{tabular}\par}
  \begin{intrebare}<2->
    Care este rezultatul $a \mathbin{\&} b$ pentru $a = 1010_2$, $b = 1100_2$?
  \end{intrebare}
  \uncover<2->{\centering
  \begin{tabular}{rcccc}
    $a$: & 1 & 0 & 1 & 0 \\
    $b$: & 1 & 1 & 0 & 0 \\
    \hline
    \at{3}{$a \mathbin{\&} b$:} & \at{3}{1} & \at{3}{0} & \at{3}{0} & \at{3}{0} \\
  \end{tabular}\par}
  \begin{uncoverenv}<4->
    \begin{lstlisting}
int a = 0b1010, b = 0b1100;
int r = a & b; // 0b1000 = 8
    \end{lstlisting}
  \end{uncoverenv}
\end{frame}

\begin{frame}[fragile]{OR \textemdash{} \texttt{|}}
  Rezultat $1$ dacă \textbf{cel puțin unul} dintre biți este $1$.

  {\centering
  \begin{tabular}{cc|c}
    \toprule
    $p$ & $q$ & $p \mid q$ \\
    \midrule
    0 & 0 & 0 \\
    0 & 1 & 1 \\
    1 & 0 & 1 \\
    1 & 1 & 1 \\
    \bottomrule
  \end{tabular}\par}
  \begin{intrebare}<2->
    Care este rezultatul $a \mid b$ pentru $a = 1010_2$, $b = 1100_2$?
  \end{intrebare}
  \uncover<2->{\centering
  \begin{tabular}{rcccc}
    $a$: & 1 & 0 & 1 & 0 \\
    $b$: & 1 & 1 & 0 & 0 \\
    \hline
    \at{3}{$a \mid b$:} & \at{3}{1} & \at{3}{1} & \at{3}{1} & \at{3}{0} \\
  \end{tabular}\par}
  \begin{uncoverenv}<4->
    \begin{lstlisting}
int a = 0b1010, b = 0b1100;
int r = a | b; // 0b1110 = 14
    \end{lstlisting}
  \end{uncoverenv}
\end{frame}

\begin{frame}[fragile]{XOR \textemdash{} \texttt{\textasciicircum}}
  Rezultat $1$ dacă biții sunt \textbf{diferiți}.

  {\centering
  \begin{tabular}{cc|c}
    \toprule
    $p$ & $q$ & $p \oplus q$ \\
    \midrule
    0 & 0 & 0 \\
    0 & 1 & 1 \\
    1 & 0 & 1 \\
    1 & 1 & 0 \\
    \bottomrule
  \end{tabular}\par}
  \smallskip
  Proprietăți XOR: $a \oplus a = 0$; \quad $a \oplus 0 = a$; \quad XOR este comutativ și asociativ.
  \begin{intrebare}<2->
    Care este rezultatul $a \oplus b$ pentru $a = 1010_2$, $b = 1100_2$?
  \end{intrebare}
  \uncover<2->{\centering
  \begin{tabular}{rcccc}
    $a$: & 1 & 0 & 1 & 0 \\
    $b$: & 1 & 1 & 0 & 0 \\
    \hline
    \at{3}{$a \oplus b$:} & \at{3}{0} & \at{3}{1} & \at{3}{1} & \at{3}{0} \\
  \end{tabular}\par}
  \begin{uncoverenv}<4->
    \begin{lstlisting}
int a = 0b1010, b = 0b1100;
int r = a ^ b; // 0b0110 = 6
    \end{lstlisting}
  \end{uncoverenv}
\end{frame}

\begin{frame}[fragile]{Deplasări pe biți \textemdash{} left shift și right shift}
  \begin{intrebare}
    Ce efect au deplasările la stânga și la dreapta asupra valorii unui număr?
  \end{intrebare}
  \begin{itemize}
    \item<2-> \lstinline|x << k| $= x \cdot 2^k$ \quad (deplasare la stânga)
    \item<2-> \lstinline|x >> k| $= \lfloor x / 2^k \rfloor$ \quad (deplasare la dreapta)
  \end{itemize}
  \begin{uncoverenv}<3->
    \begin{lstlisting}
int x = 1;      // 0b00000001
int a = x << 3; // 0b00001000 = 8   (x * 2^3)
int b = x >> 0; // 0b00000001 = 1

int y = 12;     // 0b00001100
int c = y << 1; // 0b00011000 = 24  (y * 2)
int d = y >> 1; // 0b00000110 = 6   (y / 2)
int e = y >> 3; // 0b00000001 = 1   (y / 8)
    \end{lstlisting}
  \end{uncoverenv}
  \begin{remarca}<4->{Atenție}
    \lstinline|1 << k| dă overflow pentru $k \geq 31$ pe \lstinline|int|; pentru
    \lstinline|long long| se folosește \lstinline|1LL << k|.
  \end{remarca}
\end{frame}

% ================================================================
\section{Aplicații}

\begin{frame}[fragile]{Verificarea unui bit}
  \begin{intrebare}
    Cum se verifică dacă al $k$-lea bit al lui $x$ este setat?
  \end{intrebare}
  \begin{hint}<2->
    Se aduce bitul $k$ pe poziția $0$ prin deplasare, apoi se izolează cu o mască.
  \end{hint}
  \begin{uncoverenv}<3->
    \begin{lstlisting}
// Verificam daca al k-lea bit este setat
bool isSet(int x, int k) {
    return (x >> k) & 1;
}

// Caz special: x impar <=> bitul 0 setat
bool isOdd(int x) {
    return x & 1;
}
    \end{lstlisting}
  \end{uncoverenv}
  \begin{itemize}
    \item<4-> Exemplu: $x = 13 = 1101_2$, $k = 2$: $x \gg 2 = 0011_2$, iar $0011_2 \mathbin{\&} 1 = 1$ $\Rightarrow$ bitul $2$ este setat
    \item<4-> Pentru $k = 1$: $x \gg 1 = 0110_2$, iar $0110_2 \mathbin{\&} 1 = 0$ $\Rightarrow$ bitul $1$ nu este setat
  \end{itemize}
\end{frame}

\begin{frame}[fragile]{Setarea, resetarea și inversarea unui bit}
  \begin{intrebare}
    Cum se setează, se resetează și se inversează al $k$-lea bit al lui $x$?
  \end{intrebare}
  \begin{itemize}
    \item<2-> Masca \lstinline|1 << k| are setat doar bitul $k$
    \item<2-> \lstinline+|+ setează, \lstinline|& ~| resetează, \lstinline|^| inversează bitul selectat de mască
  \end{itemize}
  \begin{uncoverenv}<3->
    \begin{lstlisting}
// Setam al k-lea bit (il facem 1)
int setBit(int x, int k) {
    return x | (1 << k);
}

// Resetam al k-lea bit (il facem 0)
int clearBit(int x, int k) {
    return x & ~(1 << k);
}

// Flippuim al k-lea bit (0->1, 1->0)
int flipBit(int x, int k) {
    return x ^ (1 << k);
}
    \end{lstlisting}
  \end{uncoverenv}
\end{frame}

\begin{frame}[fragile]{Ultimii $k$ biți și puterile lui 2}
  \begin{lstlisting}
// Ultimii k biti ai lui x  (ex: k=3 => masca 0b0111)
int lastKBits(int x, int k) {
    return x & ((1 << k) - 1);
}
  \end{lstlisting}
  \begin{intrebare}
    Cum se verifică dacă $n$ este putere a lui $2$?
  \end{intrebare}
  \begin{hint}<2->
    Ce efect are $n \mathbin{\&} (n-1)$ asupra celui mai din dreapta bit de $1$?
  \end{hint}
  \begin{itemize}
    \item<3-> O putere a lui $2$ are exact un singur bit de $1$
    \item<3-> $n \mathbin{\&} (n-1)$ elimină cel mai din dreapta bit de $1$
    \item<3-> Exemplu: $4 = 100_2$,\ $3 = 011_2$,\ $4 \mathbin{\&} 3 = 0$ \checkmark
  \end{itemize}
  \begin{uncoverenv}<4->
    \begin{lstlisting}
bool isPow2(int n) {
    return n > 0 && (n & (n - 1)) == 0;
}
    \end{lstlisting}
  \end{uncoverenv}
\end{frame}

% ================================================================
\section{Funcții utile}

\begin{frame}[fragile]{Popcount \textemdash{} numărul de biți de 1}
  \begin{intrebare}
    Cum se determină numărul de biți de $1$ ai unui număr?
  \end{intrebare}
  \begin{uncoverenv}<2->
    \begin{lstlisting}
// Implementare manuala: O(nr. biti)
int popcount(int x) {
    int cnt = 0;
    while (x > 0) { cnt += x & 1; x >>= 1; }
    return cnt;
}
    \end{lstlisting}
  \end{uncoverenv}
  \begin{uncoverenv}<3->
    \begin{lstlisting}
// Built-in GCC (O(1), foarte rapid)
int cnt = __builtin_popcount(x);    // pentru int
int cnt = __builtin_popcountll(x);  // pentru long long
    \end{lstlisting}
  \end{uncoverenv}
  \begin{itemize}
    \item<4-> Exemplu: \lstinline|popcount(13)| = \lstinline|popcount(0b1101)| = 3
    \item<4-> Built-in-urile \lstinline|__builtin_*| sunt specifice GCC \textemdash{} funcționează
      pe majoritatea compilatoarelor de concurs (g++), dar nu sunt garantate pe MSVC
  \end{itemize}
\end{frame}

\begin{frame}[fragile]{Leading zeros și trailing zeros}
  \begin{itemize}
    \item \textbf{Leading zeros}: biți de $0$ \textbf{de la stânga} până la primul bit de $1$
    \item \textbf{Trailing zeros}: biți de $0$ \textbf{de la dreapta} până la primul bit de $1$
  \end{itemize}
  \begin{intrebare}
    Câți biți de $0$ sunt la stânga, respectiv la dreapta primului bit de $1$ pentru
    $x = 12 = 00001100_2$ (pe $32$ de biți)?
  \end{intrebare}
  \uncover<2->{Leading zeros $= 28$ (primii $28$ de biți sunt $0$), trailing zeros $= 2$.}
  \begin{uncoverenv}<3->
    \begin{lstlisting}
int lz = __builtin_clz(x);  // leading zeros  (int)
int tz = __builtin_ctz(x);  // trailing zeros (int)
// Variante pentru long long: __builtin_clzll, __builtin_ctzll

// Atentie: __builtin_clz(0) si __builtin_ctz(0) sunt nedefinite!
    \end{lstlisting}
  \end{uncoverenv}
  \begin{itemize}
    \item<4-> Trailing zeros al lui $x$ = indicele celui mai mic bit setat al lui $x$
    \item<4-> Exemplu: \lstinline|__builtin_ctz(12)| = \lstinline|__builtin_ctz(0b1100)| = 2
  \end{itemize}
\end{frame}

\begin{frame}[fragile]{Lowest set bit \textemdash{} \lstinline|x & (-x)|}
  \begin{intrebare}
    Cum se izolează cel mai din dreapta bit de $1$ (\textbf{lowest set bit}) al lui $x$?
  \end{intrebare}
  \begin{hint}<2->
    Reprezentarea \textbf{complement față de 2}: $-x = \,\sim\! x + 1$ (se inversează toți biții și se adaugă $1$).
  \end{hint}
  \begin{uncoverenv}<3->
    \begin{lstlisting}
int lowestBit(int x) { return x & (-x); }
// x = 12 = 0b1100  =>  lowestBit(12) = 0b0100 = 4
    \end{lstlisting}
  \end{uncoverenv}
  \begin{itemize}
    \item<4-> La adunarea lui $1$, carry-ul se propagă exact până la primul bit de $1$ din $x$
    \item<4-> Astfel $-x$ coincide cu $x$ începând de la primul bit de $1$ (inclusiv), iar restul este inversat
    \item<4-> $x \mathbin{\&} (-x)$ păstrează doar acel bit comun
  \end{itemize}
  \uncover<5->{
    Exemplu: $x = 12 = \mathbf{1100}_2$,\quad $-x = \mathbf{0100}_2$ (în 4 biți),\quad
    $x \mathbin{\&} (-x) = 0100_2 = 4$
  }
\end{frame}

% ================================================================
\section{Recapitulare}

\begin{frame}[fragile]{Recapitulare \textemdash{} Operații uzuale}
  \begin{center}
  \begin{tabular}{ll}
    \toprule
    Operație & Cod C++ \\
    \midrule
    Bitul $k$ e setat?   & \lstinline|(x >> k) & 1| \\
    $x$ e impar?         & \lstinline|x & 1| \\
    Setează bitul $k$    & \lstinline+x | (1 << k)+ \\
    Resetează bitul $k$  & \lstinline|x & ~(1 << k)| \\
    Flip bitul $k$       & \lstinline|x ^ (1 << k)| \\
    Ultimii $k$ biți     & \lstinline|x & ((1 << k) - 1)| \\
    $n$ e putere de 2?   & \lstinline|n > 0 && !(n & (n-1))| \\
    Lowest set bit       & \lstinline|x & (-x)| \\
    Număr de biți de 1   & \lstinline|__builtin_popcount(x)| \\
    Trailing zeros       & \lstinline|__builtin_ctz(x)| \\
    \bottomrule
  \end{tabular}
  \end{center}
\end{frame}

\begin{frame}{Probleme de antrenament}
  \begin{block}{Ușoare}
    \begin{enumerate}
      \item \href{https://www.pbinfo.ro/probleme/2577/getbit}{\textbf{getbit}}
            \textcolor{gray}{(pbinfo.ro)}
      \item \href{https://www.pbinfo.ro/probleme/2580/setbit0}{\textbf{setbit0}}
            \textcolor{gray}{(pbinfo.ro)}
      \item \href{https://www.pbinfo.ro/probleme/2581/setbit1}{\textbf{setbit1}}
            \textcolor{gray}{(pbinfo.ro)}
      \item \href{https://www.pbinfo.ro/probleme/2711/puterebiti}{\textbf{puterebiti}}
            \textcolor{gray}{(pbinfo.ro)}
    \end{enumerate}
  \end{block}
  \begin{block}{Medii}
    \begin{enumerate}
      \item \href{https://code.algolymp.com/problems/4271}{\textbf{Base Bit Report}}
            \textcolor{gray}{(AlgoCode \#4271)}
      \item \href{https://www.pbinfo.ro/probleme/2799/addone}{\textbf{addone}}
            \textcolor{gray}{(pbinfo.ro)}
      \item \href{https://www.pbinfo.ro/probleme/2624/invbits}{\textbf{invbits}}
            \textcolor{gray}{(pbinfo.ro)}
    \end{enumerate}
  \end{block}
\end{frame}

\begin{frame}{Resurse}
  \begin{enumerate}
    \item \href{https://www.pbinfo.ro/articole/5562/baze-de-numeratie}{\textbf{pbinfo.ro}}
          \textemdash{} Baze de numerație (RO)
    \item \href{https://www.pbinfo.ro/articole/62/operatii-pe-biti}{\textbf{pbinfo.ro}}
          \textemdash{} Operații pe biți (RO, ușor de înțeles)
    \item \href{https://cp-algorithms.com/algebra/bit-manipulation.html}{\textbf{cp-algorithms.com}}
          \textemdash{} Bit manipulation (EN)
    \item \href{https://cppi.sync.ro/materia/lectii/baze\%20de\%20numeratie,\%20scrierea\%20constantelor\%20si\%20operatii\%20pe\%20biti.pdf}{\textbf{cppi.sync.ro}}
          \textemdash{} Material mai stufos de Marius Nicoli (RO)
    \item \href{https://en.wikipedia.org/wiki/Two\%27s_complement}{\textbf{wikipedia.org}}
          \textemdash{} Two's complement (EN)
  \end{enumerate}
\end{frame}

\end{document}
